Ad aliud, de canone librorum, ubi nixus Hieronymi et concilii Florentini auctoritate mavult credere concilio, quod a spiritu sancto regitur, quam mihi, et gratias ago. Pie enim sapit, nunquam volui mihi credi. Sed respondeo breviter: Conciliet ipse primo Hieronymum sibi, qui in prologo galeato Machabaeorum libros et nonnullos alios manifeste inter apocrypha recenset, qua auctoritate fit, ut mihi liber Machabaeorum sit gratus et probatus, sed contentiosis pateat ad repulsam. Satis hodie de concilio dixi, imo, ut ipsemet D. D. dicit: Concilium non errat, si autem errat, non est concilium, et ut meo sensu loquar, credo concilium et ecclesiam nunquam errare in his, quae sunt fidei, in ceteris non est necesse non errare.

Quod vero exclamat, latibulum hoc esse haereticorum, qui hac fiducia quaslibet expositiones auctoritatum refutarent vel acciperent, respondeo: Quare hoc non sunt conquesti s. Augustinus, Hieronymus aliique victoriosissimi haereticorum triumphatores, sed donatis ambiguis locis studuerunt certis et apertis locis pugnare? Hoc faciamus et de animabus in purgatorio.

Ad rem et auctoritates.

Ad solutiones auctoritatum dixit, a me contemptos esse Augustinum, Hieronymum, Ambrosium, Bernhardum, Gregorium, Isidorum cum glossa ordinaria. Hoc pro modestia, imo molestia Ecciana dixit, nimis cupidus movendae invidiae. Dixi hodie, non esse contemptos a me, et iterum dico.

Quod autem meam glossam non accipit, non curo. Probet autem ipse et suam: cum textus sit clarus de die Domini et igne, in quo revelabitur dies Domini. Quod etsi potest trahi ad purgatorium, ut dixi et sic sententiam D. doctoris \marginpar{\small[340]} non reprobavi, purgatorium scilicet esse, non tamen potest hac intelligentia oppilari os contentiosum.

Item et illud Matth. 5. [v. 25] de carcere non contempsi pro Ambrosii sententia. Sed quod addit D. D. verba Christi indicare, debere solvi usque ad novissimum quadrantem, ideo de inferno intelligi non debere: hoc satis eluitur per divum Hieronymum, qui adversus

Helvidium hanc dictionem »donec« recte exponit, quod non cogatur significare, quod Ambrosius cum D. D. hic voluit. Nam sic et Matth. 1. [v. 25] scribitur: Nec cognovit eam, donec peperit filium suum, et tamen non sequitur, quod post partum cognoverit eam, ita hic non sequitur, ut sit exiturus post solutionem, sed solvet et non exibit.

Quod autem confutavit, quod non peccaverit mortaliter dissentiens adversario, atque quod et mortalia, modo sint contrita, et venialia purgentur in purgatorio: dico, quod hic textus non loquitur de contritis mortalibus, nam contritum mortale iam non est mortale, et poenitens iam non est dissentiens adversario suo, ergo de dissentiente et sic permanente loquitur Christus.

Transeo illud de casu ligni apud Ecclesiasten, ubi Hieronymum inducit intelligere per austrum statum bonum et per aquilonem statum malum, concedo utrumque, et nihil ad purgatorium.

De mansione apud Iohannem certa, ubi argutatur in hunc modum: Quomodo fiet certa, si fiat accessio gratiae, cui debeatur altior mansio? potest idem argui de quolibet fideli post baptismum, cui est certa mansio ab aeterno praedestinata.

Deinde in purgatorio existentes sic ordinati sunt, ut hoc modo ad certam mansionem veniant. Illud etiam humanum argumentum nihil concludit, quod sequeretur descendentem cum pluribus venialibus peccatis melius habiturum esse quam eum, qui cum paucis. Si haec ratio bona est, timendum est, ne meretrix melior fiat quam b. virgo, quasi non sint differentes gradus animarum in purgatorio, cum ipse hodie sanctos velut prae ceteris praestantiores in purgatorio asseruerit.

Deinde Silvestrinum argumentum inducit, quod prodesset eis, si perdurarent in poenis, nam sic et martyribus expediret usque in diem iudicii mori et pati, quasi ignoret D. D. ad certam mensuram poenas esse animabus constitutas. Non enim damnabile est, orare pro mortuis, sicut nec damnabile fuit, quod apostolus pro se oravit et orari petiit, cum tamen cresceret virtus eius in infirmitate. Ita quilibet fidelis pro qualibet necessitate cuiuslibet fidelis debet orare, debet adiuvare, non obstante, quod ille per necessitatem istam magis ac magis mereatur.

Per idem volo solutum, quod contra observantiam ecclesiae et pietatem sit, quod dixi animam alienis suffragiis liberatam minus beari, quam si per \marginpar{\small[341]} se satis patiatur in purgatorio. Hoc velut corollarium sumpsit ex dictis meis, ubi ego sensi animabus esse succurrendum, quantumlibet perfectius genus meriti esset iustitiae divinae omnibus modis satisfacere.

Damascenum transeo cum termino meriti et remitto ad praecedentia.

Item, iterum me carpit, quod nullo fulcitus testimonio prorogo terminum meriti, ipse autem nitatur multis et neget. Respondeo: Nullis nititur et ipse nisi violenter intellectis, ut satis dictum est.

Capere etiam me in verbis meis volens dicit, me appellasse sententiam meam opinionem, ergo male oppositum errorem appellaverim. Dico sicut prius. Non solum opinionem, sed et ignorantiam appello meam sententiam, errorem appellavi non contrariam opinionem, sed quod opinionem pro veritate statuunt.

Deinde auctoritatem Hieronymi de Noe, Job et Daniele dicit valere non tantum pro extremo iudicio, nec possim hoc convincere ex litera. Relinquo hoc iudicio melius sentientis, cum manifeste textus habeat, quod pro homine ad tribunal, quod in morte fit, rapto nec Noe nec Daniel nec Job orent, quod omnino de iudicio extremo intelligi oportet, aut saltem non de purgatorio, quia pro mortuis in purgatorio orant Noe, Job, Daniel et omnis ecclesia.

Ad auctoritatem Apocal. recte dicit, quod daemones sub terra 'non cantant, sed blasphemant etc.'.

D. D. Martinus exclusus tempore, cum iam duae horae transissent, obtulit se ad reliqua in scriptis tradenda notariis et primum D. doctori Eccio exhibenda. Obtulit autem infrascripta per notarium suum.

Quod daemones in inferno non laudant, sed blasphemant deum, scio, sed quod ideo subtus terram significet purgatorium, in quo cantent animae, non sequitur. Primum, quod sunt in tribulationibus et poenis necdum abstersit deus omnem lacrimam ab oculis earum. D. D. autem negat, cantare esse in tribulationibus laudare deum. Ideo

contra se ipsum loquitur, dum subtus terram in purgatorio laudem ponit simul et poenas, in quibus non laudent. Nego autem et ego, quod animae ideo cantent, quia habent certissimum signum beatitudinis. Hoc enim signum certissimum fuit probandum. Deinde ego sub terram non dixi significare solum infernum, sed et alia, quaecunque sub terra sunt, in quibus omnibus laus dei abundat, omnia enim sunt plena gloria et laude dei, qui laudabilis est in omnibus operibus suis. Alioquin oportet D. doctorem non solum purgatorium per subtus terram intelligere, sed etiam aliquod quartum per subtus mare, cum eodem Apocal. 5. etiam subtus mare recitetur et omnia, quae in eo sunt. Rectius ergo creaturas omnes hoc quaternario intelligimus significatas, quam ea, quae D. doctor Eccius imaginatur.

\marginpar{\small[342]} Quod autem in tribulationibus laudare deum sit cantare, probo non esse a me fictum, ut D. D. dicit, quia Rom. 5. [v. 3.] dicit: Gloriamur in tribulationibus, et Psal. 42 [v. 9]: In die mandavit Dominus misericordiam suam et nocte (id est, tribulatione), canticum eius, et Jac. 5 [v. 13]: Tristatur aliquis, oret aequo animo et psallat; item: Benedicam Dominum in omni tempore [Ps. 34 v. 2]. Omni, inquit tempore, etiam adverso. Contra vituperatur ille, de quo Psal. [49 v. 19]: Confitebitur tibi, cum benefeceris ei. Fateor autem hoc me non dixisse de damnatis subtus terram, sed ut excluderem sententiam D. doctoris volentis, quod cantare sit tantum de beatitudinis securitate in purgatorio.

De canone missae non est contentus, quod dormire in pace ad corpus retuli, quia dicit canon: In Christo quiescentibus, non in sepulcro quiescentibus. Porro quid referat ad animam, corpus in mari aut aëre versari ac non quiescere? Item, quis (inquit) corpus exanime dicat dormire in pace? Respondeo: Nihil ista argutia valet, nam quiescentes in sepulcro aut ubilibet secundum corpus iactati vere quiescunt in Christo, hoc est ut Apocal. 14 [v. 13]: requiescunt a laboribus suis, videlicet quod separata anima a corpore iam non laborat in corpore in variis huius mundi molestiis. Alioquin quomodo Christi corpus in sepulcro negabit dormire in pace et requiescere in deo, cum in tot locis de eo sic loquatur scriptura et cantet ecclesia: In pace in id ipsum dormiam et requiescam [Ps. 4 v. 9]; et Psal. 16 [v. 9]: Et caro mea requiescet in spe. Item Apocal. 14 [v. 13]: Beati qui moriuntur in Domino. Item Stephanus Act. 7 [v. 59]: Obdormivit in Domino, quod sine dubio de corporali morte tropo scripturae dicitur. Si Eccio

non dormit corpus exanime, sinat, quaeso, illud dormire suscitatori deo et scripturae auctoritati.

In fine dicit, orari a nobis pro pace, non quam habent, sed quam habere a miseria poenarum contingat purgatis. Respondeo: Haec est petitio principii, quia pax illa securitatis, quam habeant, nondum est probata. Relinquitur ergo, quod in poena sunt, et ut ab hac solvantur et pacem habeant, oremus.

Haec omnia sic volo dixisse, ut ostendam me ignarum esse eorum, quae in purgatorio agat deus, et nimio eos audere, qui huius ignorantiae impatientes potius fingere volunt sua, quam confiteri sese ignorare. Mea probare non possum, sed nec illi sua. Quare recte mea conclusio id tantum voluit, esse nec scripturis nec rationibus probatum, animas certas esse saltem omnes de salute sua et gratiam in eis non augeri, sufficiat, quod scimus, eas pati et nos debere eis succurrere, cetera deo relinquenda soli.

Acta sunt haec in praesentia honorabilium virorum Johannis Teuber Weysmann et Bartholomaei Schaller Ernfriedsdorfensis, Bambergensis et Misnensis dioecesis, laicorum testium ad praemissa requisitorum specialiter et vocatorum. Die solis decima mensis Julii in solito disputationis loco, hora meridiana vel quasi.

D. D. Eccius sequentia obtulit in scheda quadam Die lunae undecima mensis Julii mane hora septima.

Solutiones reverendi patris apparent mihi insufficientes, potissimum quod Augustinum dicat non loqui de purgatorio, cum maxime hoc agat, scilicet \marginpar{\small[343]} de purgatorio. Neque negandum est, animas defunctorum pietate suorum viventium relevari, cum pro illis sacrificium mediatoris offertur vel eleemosynae in ecclesia fiunt, sed haec eis prosunt, qui, cum viverent, ut haec sibi postea prodesse possent, meruerunt. Est enim vivendi modus nec tam bonus, ut non requirat illa post mortem, nec tam malus, ut ei non prosint post mortem. Quocirca hic omne meritum comparatur, quo possit post hanc vitam relevari quispiam vel gravari. I nunc, Luthere, et dic Augustinum hic non de purgatorio locutum.

Praeterea non probat, ad deletionem peccati venialis requiri augmentum gratiae. Fateor, augmentum gratiae tollere aliquando veniale, sed etiam veniale tollitur satispassione. Unde apostolus allegatus Rom. 7. [v. 14 ff.] non loquitur de veniali peccato, porro si loqueretur, satispassionem non excluderet. At de eius consensu, velut impertinente, nihil dico in praesentia. Veniale, fateor, offendit, sed venialiter, quia hominem in odio dei non constituit nec inimicum dei facit, quod alioquin cum gratia non staret, ideo in purgatorio quidem habent peccata venialia citra inimicitiam dei.

At quod causatur, me torquere auctoritates per capillos ad propositum, dispeream, si tota disputatione vel unam adduxerit auctoritatem tam pertinentem, sicut ego hic adduxi. Sed hoc iudicent iudices, non Lutherus.

Indignatur, cur non et Augustinus et Hieronymus causati sint hoc latibulum haereticorum, quod dixi; bone deus, si non ubique clamet Augustinus contra Pelagianos, Cyprianus et Ambrosius contra Arianos et Novatianos, eos pro arbitrio interpretari sacram scripturam!

Petit, quod probem meam sententiam, cum suam probare non possit, meam vero probatam iudicum relinquo iudicio.

De »donec« alia est ratio Hieronymi ad Helvidium, quia aliunde probatur ex scriptura, Mariam semper mansisse virginem, igitur donec ibidem non dicit consummationem, quod hic non potest facere.

Porro de rationibus meis gratia brevitatis non moror. Sed quod dixit, me certum fecisse ex dictis suis, quod animae suffragiis viventium liberatae minus beentur, male mihi imponit, nec enim mihi somnia fingo. Refert enim in dilutione adnotationum mearum, noluisse Paschalem sua sibi minui praemia, ideo maluisse ardere, sed missa haec facio et iudicibus committo.

Porro, non nego, quin in tribulationibus quis laudet deum, sed hoc loco accipi non potest, quod et in caelo cantare eos dicit, cum ab illis absterserit deus omnem lacrimam. Unde liquet lectori, quomodo sententiam meam invertat. Unde animae in purgatorio cantant, secundum Johannem, deo, et quaero, causam explicet incertitudinis gloriae; quare falsissime comminiscitur, quod contra me ipsum loquar. Unde non nego, bonos gloriari in tribulationibus et cantare, sed negavi illam glossam, quam ipse finxit, quod \marginpar{\small[344]} cantare sit in tribulationibus deum laudare, alioquin angeli non canerent deo et beati.

Quod errorem patentissimum defendere conatur, scilicet animas dormire somno pacis quoad corpus etc., quoniam anima separata a corpore

non laboret in corpore variis molestiis: vide vafritiem, quia cum solvendo ad corpus retulisset in sepulcro quiescens, iam se transfert ad quietem animae a laboribus corporis. Sed heus bone vir, est hoc quiescere a laboribus, quia videlicet non in corpore molestias patiatur, tamen incomparabiliter maiores molestias sentiat extra corpus in anima? quid iuvat, si quiesco a febri et infestor calculo?

De Christi quiete aliud est, resurrectionem absque corruptione expectante. Beati ergo, qui moriuntur in Domino. Sed quis dicturus est corpus beatum nisi per redundantiam animae in corpus? Sic et de b. Stephano, qui in Domino obdormivit.

Deinde inquit me petere principium: Sunt in pace, et tamen eis petimus pacem. Et quia ridiculum est, referre pacem habitam ad corpus, sequitur, illam esse in anima scientiam securitatis.

Quod ait, se sua probare non posse, reputo verissimum, contra alii sua probant, quamvis pertinax et cervicosus etiam demonstrationibus non acquiescat, quoniam semper contenderet, non esse demonstrationem, velut Gregorius Arimin. et Petrus Aliacensis Aristoteli faciunt in rationibus de primo motore in 7. et 8. Physicorum. Quare cum communi concludo sententia, et in purgatorio gratiam non augeri, et eas certas esse de salute.

Die lunae undecima mensis Julii hora septima coepit disputari materia de indulgentiis D. D. Eccius.

Imprimis, antequam descendam in hanc disputationem indulgentiarum, protestor, non esse mentis meae aut propositi velle hac mea disputatione contravenire mandatis summi pontificis, quibus mandat in decreto incipiente: Cum postquam, non debere defendi vel praedicari contrarium certis punctis de indulgentiis sub poena excommunicationis, sed magis esse intentionis meae, defendere veritates in eodem decreto approbatas. Quo sic praefato, accedo ad impugnandum conclusionem undecimam.

Dicere indulgentias esse utiles Christianis, est verum et pium, nec indulgentiae sunt vitium boni operis, quare non videtur bene dictum, quod \marginpar{\small[345]} hoc affirmantes insaniant. Hoc probatur, quia ecclesia in iis, quae sunt fidei, non errat, et quae respiciunt animarum salutem. Imo, ut s. Cyprianus ait, deus non permittit maiorem partem cleri errare, sed in conciliis generalibus et per universitatem cleri iam a trecentis annis indulgentiae reputatae sunt utiles et piae Christianis. Quod patuit in concilio Viennensi, ubi approbatae sunt indulgentiae datae per Urbanum IV. pro venerabilis eucharistiae veneratione, nam concilium hanc addidit rationem motivam, ut Christi fideles essent magis parati ad obsequia venerationis et honoris sacramento impendenda.

Sic Innocentius III. in concilio generali, ubi utilissimam constitutionem: Omnis utriusque a tota ecclesia receptam edidit, de indulgentiis quoque disposuit hospitaliorum, dicens: Iniungimus vobis in remissionem peccatorum, quatenus de bonis vobis collatis grata eis charitatis subsidia erogetis, ut per subventionem vestram ipsorum inopiae consulatur et vos per haec bona et alia, quae Domino inspirante feceritis, ad aeterna possitis gaudia pervenire. In eodem quoque concilio limitata fuit inferiorum praelatorum indulgentias dandi potestas, quae in Lugdunensi concilio deinde fuit approbata. Sed si indulgentiae essent vitium boni operis et Christianis inutiles, quid opus esset collationem vitii operis minuere in inferiori praelato et conferre in plenitudinem potestatis?

Accedit, quod laudabile Constantiense concilium, quod et inter errores damnavit indulgentiarum contemptum, contulit indulgentias omnibus in concilio existentibus et feria sexta ieiunantibus, quod fuisset bona opera tantorum patrum inquinare et inficere. Et Gerson, quem reverendus pater illustrem theologum appellat, sicut fuit veri et honesti studiosissimus, is decidit, indulgentiarum concessionem non esse parvi pendendam vel contemnendam, sed devote amplectendam in fide, spe et charitate Domini nostri Jesu Christi, qui talem potestatem clavium ecclesiasticarum dedit hominibus. Constat enim, quod fructuosior est et deo acceptabilior operatio talibus indulgentiis innitens, quam altera et hoc ceteris paribus non innitens. Haec Gerson.

Sic et alii sancti et optimae existimationis viri tempore datarum indulgentiarum voluerunt, quae a b. Gregorio quoque datae sunt ante

nongentos annos et Paschasio ante sexcentos annos. Unde sola tanti patris auctoritas debet quemlibet Christianum movere, ut credat indulgentias esse utiles Christianis, quamvis non ignorem, quod Vuilhelmus Altisiodorensis referat non defuisse, qui etiam vivente b. Gregorio sancto patri in indulgentiis contradixerint.

Porro hoc addendum: Si indulgentiae sunt vitium boni operis, hoc potissimum meo arbitratu ex eo suspicari posset, quod sunt satisfactoriae, sed \marginpar{\small[346]} per opus satisfactorium non minus meremur, quam si non esset satisfactorium, alioquin praestaret nos nulla facere opera satisfactoria, ne merita nostra pro isto statu minueremus. Sed in benedictionibus seminantes de benedictionibus quoque meteremus. Sed hic audiam reverendum patrem. Accedat tantus totius Christianitatis consensus in iubilaeis per summos pontifices celebratis, per Bonifacium, Clementem VI., Urbanum VI., Nicolaum V., integerrimum et doctissimum pontificem, Sixtum IV., et quod communi consensu Christianissimorum regum et principum pro passagiis et cruciatis summus pontifex plenissimas saepe dedit indulgentias.

Cum ergo fides ecclesiae sit indefectibilis, pro qua Christus rogavit, ut non deficeret, Lucae 22, et pollicitus est, se nobiscum esse omnibus diebus usque ad consummationem saeculi Matth. 28. [v. 20], non est admittendum, ecclesiam cum animarum pernicie tanto tempore errasse, et quia decreta summorum pontificum ab omnibus communionem ecclesiae sanctae habentibus sunt acceptanda, ut ait Gregorius can. Praeceptis, 12. distin., item can. Omnia decretalia, 25. quaest. 1. cum multis concordantibus, et pontifex modernus sub poena excommunicationis latae sententiae definivit, indulgentias esse utiles ad poenam pro peccatis debitam his, qui charitate iungente membra sunt Christi, delendam, et quod huiusmodi remissio ex merito Christi et sanctorum, auctoritate summi pontificis, quasi compensetur: his omnibus persuasum volo Christiano fideli, indulgentias non esse vitium boni operis, nec dicere indulgentias Christiano utiles esse insanire, quod nihil aliud est, quam dicere, totum clerum in ecclesia tanto tempore insanivisse.

\paragraph*{D. D. Martinus.}

Duodecim argumentis impugnat egregius D. D. conclusionem meam undecimam. Ad quae priusquam respondeo, primo diluam illud,

quod insaniam esse dixi, habere indulgentias bonas Christiano esseque eas vitium boni operis. Dico ergo, quando propheta Psal. 40. [v. 5] audet insanias falsas appellare, et Psal. 119. [v. 85] iniquorum fabulationes ipsas traditiones hominum, quae de se non erant malae, nisi quod malam fiduciam adversus iustitiam dei praestabant insipientibus: quanto magis insania est indulgentias, quae neque praeceptae neque consultae neque necessariae ad salutem, imo remissiones sunt bonorum multorum operum, haberi pro bono Christianorum? cum omne bonum nostrum et necessarium ad salutem comprehensum sit in praeceptis et consiliis. Imo amplius dico: Quando saluberrima dei lex, teste Ezechiele et apostolo Paulo Rom. 7. [v. 11 ff.], non est bona homini, imo occasio peccati, sola autem gratia bonum est Christiano, quanto minus indulgentiae, quae nedum cum gratiae, sed nec cum legis ullius bonitate ullo modo comparari possunt, \marginpar{\small[347]} cum sint nihil et privatio quaedam bonorum operum, haberi debent pro bono? Non quod noxias aut perniciosas arbitrer, ut dicemus.

Quando ergo egregius D. D. dicit primo, ecclesiam 'in his', quae sunt fidei, non errare, nec in his, quae respiciunt animarum salutem, recte dicit, sed tales non sunt indulgentiae, ut dictum est.

Quod autem Cyprianus dicit, deum non permittere, ut maior pars cleri erret, viderit ipse, quid dixerit. Certum est, quod tempore Arianae perfidiae maior pars et eloquentiorum episcoporum et doctiorum virorum erravit in tanto articulo fidei, adeo ut nullus catholicus episcopus ferme resideret in cathedra sua.

Quod dicit, in conciliis generalibus et a trecentis annis habitas indulgentias utiles et pias Christianis et in concilio Viennensi approbatas, quas Urbanus IV. pro reverentia venerabilis eucharistiae in fidelibus excitanda largitus est, respondeo: Nunquam negavi esse utiles indulgentias, sed non Christianis, hoc est, fervide agentibus, Christum secundum nomen suum amantibus et quaerentibus, his enim ingrata est remissio operum, grata vero impositio operum. Secundo dico, quod in his rebus non est respicienda persona ullius sive concilii sive papae, sed quid dicatur, praesertim quando hic error de indulgentiis, ut in re non necessaria, non est periculosus, ceteris tamen paribus, hoc est, si sciant esse non necessarias nec in eas confidant.

Ad secundum, de concilio generali et Innocentio III., qui iniunxit pro dispositione hospitaliorum, ut in remissionem peccatorum conferrent

grata subsidia etc., dico me adhuc non intelligere hoc verbum papae, cum remissio peccatorum mihi videatur esse aliud quam indulgentiae, atque si esset idem, dico sicut prius.

Ad tertium, de Lugdunensi concilio, in quo approbata dicitur et limitatio potestatis conferendi indulgentias inferiorum praelatorum, admitto totum. Quod autem sequitur: Si essent inutiles, non fuisse necessarium collationem vitii operis minuere in inferioribus praelatis; dico indulgentias esse vocatas a me vitium operis, secundum quod impugnatus sum a D. D. conclusione mea 42. in Resolutorio, ubi sic dixi: Non quod veniae sint malae et noxiae, sed quod abusus perversus nocet, dum tale opus non facerent, nisi veniae essent. Sic enim finis operis huiusmodi fit ipsa venia, et clarissimum est, quod raro homines contribuerent, nisi indulgentiae promitterentur. Ideo semper est ibi periculum, saltem vitiosi operis.

Quarto induxit laudabile Constantiense concilium, quod inter errores ceteros damnavit contemptum indulgentiarum. Respondeo: Nunquam ego contempsi aut contemnendas esse docui, nisi contemptus intelligatur hoc modo, quod incomparabiliter meliora possimus eisdem expensis facere, quibus indul-\marginpar{\small[348]}gentiae redimuntur, vel quam sunt indulgentiae ipsae. Quare praerogativa melioris non est contemptus deterioris, sicut aurum non est contemptus ligni aut foeni, etsi praestantius.

Per idem ad quintum, ubi Gersonem inducit, indulgentiarum concessionem non esse parvi pendendam, addo, in suo genere, seu sicut ipse dixit, ceteris paribus. Nec credo hanc sententiam D. doctori satis placere, quod indulgentiae sint suscipiendae in fide, spe et charitate, cum tenuerit satisfactionem impleri et remitti posse etiam his, qui sunt extra gratiam. Admitterem, quod fructuosior sit operatio innitens indulgentiis quam altera non innitens ceteris paribus, modo cetera illa paria recte comparentur, sicut non nego, meliorem esse posse unam orationem dominicam laici, quam omnes horas canonicas unius sacerdotis.

Sexto inducit b. Gregorium ante nongentos annos, Paschasium ante sexcentos annos dedisse indulgentias, quamvis, recitante Altisiodorensi, aliqui contradixerint s. Gregorio. Respondeo: Nondum habeo fide dignam historiam de indulgentiis Gregorianis. Tamen quidquid sit, non ideo sequitur, indulgentias esse aliud quam indulgentias, hoc est, remissiones bonorum operum, quas nemo bonas poterit appellare Christianis saltem, quantumlibet sancti, concilia vel quicunque eas dederint vel approbarint. Ideo D. D. non contra me agat solum nominibus auctoritatum, sed et ipsius rei veritate.

Septimo suspicatur, ideo esse eas boni operis vitium, quia sunt satisfactoriae. Dico, quod non sunt satisfactoriae, sed remissiones satisfactionis nec habendae ut opus satisfactorium. Ideo male hoc argumento comparantur operibus satisfactoriis, quasi ideo merita nostra minuantur, si indulgentiae relinquantur, imo contrarium verius est: merita augentur, dum indulgentiae relinquuntur.

Octavo, dicit consensu fidelium receptum annum iubilaeum Bonifacii VIII., Clementis VI., Urbani VI., Nicolai V., Sixti IV. etc.. Respondeo: Quis scit, si consensus fidelium eum receperit? Nec est Bonifacius tantae opinionis, ut, quidquid gesserit, statim pro re bene gesta habendum sit, praesertim in re ista non bona nec necessaria ad salutem, qui et alia monstra perpetravit, quae in rebus necessariis malum exemplum praebuerunt.

Nono, quod dederunt Romani pontifices consensu Christianissimorum regum et principum indulgentias pro passagio etc., dico sicut prius, non est miraculum, quod deus haec contemnat et sinat agi et agere, quae impertinentia sunt ad salutem, qui in hoc super nos vigilat, ut doceat utilia, ut apostolus ad Titum 3. [v. 8] dicit, hoc est, praecepta sua et consilia.

Decimo arguit, fidem ecclesiae esse indefectibilem et Christum nobiscum, ideo ecclesiam non potuisse errare tanto tempore cum animarum pernicie. \marginpar{\small[349]} Dico: non tota ecclesia erravit. Deinde si etiam erraret in his rebus nihili (semper de Christianis loquor), non est periculum salva fide. Variari opiniones et errare in temporalibus, non tollit ecclesiam Christi.

Undecimo decreta Romani pontificis acceptanda esse contendit. Respondeo: Sine dubio sunt acceptanda, sed tamen cum iudicio, ut habetur dist. 19. cap. Anastasius; quia Romani pontifices homines fuerunt et homines circum se habuerunt, peccare potuerunt. Deinde nullum est decretum, quod praecipiat indulgentias redimere, nec est usque hodie, quod indulgentiarum valorem aperte declaraverit.

Ultimo, de moderni pontificis definitione, esse videlicet utiles indulgentias et de merito Christi et sanctorum acceptas, respondeo: Nec ipsa satis exprimit nec probat vel una syllaba, quae dicit: de qua re copiosius in actis meis Augustensibus 'dixi'. Non sequitur ergo, totum clerum in ecclesia tanto tempore insanivisse, cum interim semper fuerint, quibus et displicuerint, qui et contradixerint saltem privatim indulgentiis, ut testatur vulgatissimum proverbium: Indulgentiae sunt piae deceptiones.

\paragraph*{D. D. Eccius.}

Accepi reverendi patris sententiam in multis mitiorem et aequabiliorem, quam verba conclusionum praetenderent aut alia sua scripta mihi videbantur significare. Neque mei est instituti, indiscretas indulgentiarum proclamationes aut abusus defendere, sed veritatem et quantum indulgentiae rationabiliter datae possint prodesse, tutari. Et ut sim brevis, non discutio de excusatione suae conclusionis. At cum assumit indulgentias non praeceptas, non necessarias esse, fatemur, Sed in hoc toto meo exili iudicio apparet esse difficultas nostrae dissensionis principaliter, quia reverendus pater dicit, indulgentias non esse, nisi remissiones bonorum operum, quam interpretationem apud nullum indulgentias vel dantem vel admittentem memini me reperisse, quoniam sic indulgentiae essent vere noxiae, periculosae et damnabiles, cum sic concilia et summi pontifices remitterent hic bona opera, ut vel sic homo non mereretur tantam beatitudinem, quod est contra animarum salutem, vel homines decepti remissione bonorum operum ad purgatorium mitterentur, ut poenas hic non dilutas ibidem sustinerent, quod revera damnabile est. Unde omnium vel dantium vel admittentium indulgentias unus consensus est, indulgentias non esse bonorum operum remissiones, sed remissionem esse poenae temporalis pro peccato contrito adhuc debitae et non solutae. Sicuti Leo modernus papa expresse in praefato decreto definivit, quod Sixtus papa IV. ante eum expressissime fecit tempore iubilaei sui.

Unde concordare non possum, quod reverendus pater in suis resolutionibus dicit, indulgentias solum datas pigris, quod \marginpar{\small[350]} etiam in sermone vulgari affirmat, non volentibus implere canones poenitentiales. Et tamen in eodem resolutorio dicit canones poenitentiales pigris impositos, quod viri alioquin strenui totam vitam poenitentiam agerent. Unde infelicissimae essent indulgentiae, si essent bonorum operum remissiones.

Soluturus ad primum respondit, indulgentias non attinere ad animarum salutem, cum profecto maxime intersit salutis animarum, an salutem promoveant vel impediant. Interest etiam fidei nostrae, an papa, ut Christi vicarius, poenas peccato debitas facta recompensatione per merita Christi et sanctorum possit remittere et delere.

De Cypriano obiecit mihi, quod occurrebat Cyprianum legenti, quamvis non plene mihi constet, an sicut in Graecia, ita et in aliis regionibus ecclesiae maior pars cleri fuerit infecta.

Deinde ad concilia respondit, indulgentias non esse bonas Christiano, id est, qui fide et charitate Christum sequatur. Hoc certe dicta conciliorum destruit et indulgentiarum largitionem, quia non nisi Christi fidelibus et contritis et confessis sub communi forma conceduntur. Quare communi sententia indulgentiae non sunt bonae malis, quia eas non consequuntur, sed bonis et vere Christianis.

Non accepto, quod reverend. pater dictum Senecae hic usurpat: Non attendendum, quis, sed quid dicatur, quoniam maxime in his, quae sunt fidei, attendere oportet etiam, quis dicat, et concilium legitime congregatum est ab omni Christiano audiendum.

Dicit errorem non periculosum, si hic de indulgentiis erretur. At mihi contrarium videtur, cum periculosi errores fuerint iudicati, qui sacris ordinibus vestris fratrum mendicantium ab initio detraxerunt.

Ad Innocentii in generali concilio constitutionem dicit, se non intelligere, quomodo valeant, vel indulgentiae sint remissio peccatorum. Dico, ut semper, non dabo, tam solenne et legitimum concilium potuisse errasse, cum a spiritu sancto non derelinquatur et remissio peccatorum intelligenda sit quoad poenam, quia nomen peccati non semper pro culpa, sed etiam pro poena culpae debita usurpatur, sicut apud Machabaeos, ut a peccatis solvantur, quod de culpa nequit intelligi mortali, sed de poena culpae debita. Et si indulgentiae dicantur dari, ut absolvantur a poena et culpa eas consequentes, nullus existimare debet, indulgentias remittere culpam, sed quia papa concedit, ut a culpa, mediante sacramento poenitentiae, a deputatis commissariis absolvantur et dein indulgentias, id est, poenarum remissiones consequantur.

\marginpar{\small[351]} Quod vero se excusat, quod dixit indulgentias esse vitium boni operis, quia et ipse voluerit ceteris paribus opus bonum cum indulgentiis esse melius quam sine indulgentiis, sed sicut fieri soleat, quia ita homines indulgentias redimant, ut bona opera illa sine indulgentiis non facerent.

Exactis iam duabus horis reliqua in horam secundam pomeridianam dilata sunt; qua veniente prosecutus est, ut sequitur.

Quamvis pleraque super illo dicto afferri possent, et in simili, quando vicini propter amicitiam vicinorum conveniunt ad oblationes et offertoria plus quam intuitu mortuorum, attamen mihi apparet, meliori iudicio semper salvo, hoc non esse vitium boni operis. Primum ex eo, quia sacrum concilium per indulgentias vult excitare homines ad bona opera. Secundo, quia esto, quod tale quid faciant vel exhibeant propter indulgentias consequendas, tamen ista opera ex illo fine non inficiuntur, quoniam finis iste indulgentiarum non discrepat a fine ultimo, ita quod poenarum remissio congruenter in deum ordinatur, alioquin praescindendo subordinationem finium paucissima vel nulla essent bona opera.

De contemptu indulgentiarum per Constantiense concilium improbato dicit reverendus pater, se illas non contempsisse, fateor ingenue, si eo semper fuisset usus moderamine, sicut se hodie exposuit, et forte eius interpretationi est standum et in conscientia sine 'foro', tamen scriptus eius sermo vulgaris aliam exhibet intelligentiam laicis. Sed hoc nihil ad me.

Ad Gersonem respondit admittendo, tamen arbitratur Gersonem mihi adversari, qui indulgentias tunc utiles reputat, quando fide, spe et charitate percipiuntur. Quod ego alias tenuerim posse fieri satisfactionem in peccato mortali, dico et ipsum Christianissimum cancellarium hoc tenere, aliquem posse satisfacere de iniuncta poenitentia in peccato mortali, alioquin imprudentissime agerent confessionum auditores, durantem imponendo poenitentiam. Sed an quispiam indulgentias in peccato mortali consequatur, nihil ad rem, ratio diversitatis est obvia, quoniam indulgentiae non solent concedi, nisi contritis et confessis,

ideo non mirum, si peccatores harum non sunt capaces, quia indulgentiae tantum valent, quantum sonant.

Subiunxit reverendus pater, orationem dominicam laici tantum valere, sicut horas canonicas clerici, si cetera sunt paria, illud non admitto.

De b. Gregorio dicit, se nondum receptam vidisse historiam. At tam constantissima fama, sicut b. Augustinus lib. de trin. de Alexandria inquit, hoc negari non permittit. Negat autem indulgentias esse satisfactoria, sed solum sint remissiones bonorum operum. In hoc credo principali nos dissentire et magnam partem controversiae hic fundari. At prius memini, indulgentias esse non remissiones bonorum operum, ut credit reverendus pater, \marginpar{\small[352]} sed remissiones poenarum peccatis debitarum. Unde homo vere Christianus accipiens indulgentias non minora, sed plura facit bona opera, quam si non consequeretur indulgentias. Sic autem intelligo indulgentias esse satisfactoria, sicut papa Sixtus IV., Clemens VI. et papa modernus declarant, non quod consecutus indulgentias satisfaciat, sed quia nullum malum manet impunitum et nullum peccatum manet inultum (can. Sicut primi de poenit. dist. 1.). Ita vicarius Christi dispensator eius thesauri, ubi homo pro poena peccati non reddit aequivalens, ipse ex thesauro sibi commisso merito Christi sufficientissimo illud recompensat et pro eo satisfacit.

Quare concessi in prioribus, indulgentias de se non delere culpam, sed poenam, et quod sit talis remissio poenarum in indulgentiis et non remissio bonorum operum, patet per illud Esaiae: Spiritus Domini super me, eo quod unxerit me, ad annunciandum mansuetis misit me, ut mederer contritis ac praedicarem captivis indulgentiam et clausis apertionem etc.. Christus eadem verba de se assumens pro verbo indulgentiam remissionis nomen accepit. Sic Theophilus ait, ut b. Thomas eum enarrat: nam prius ab inferis animae omnium tenebantur, donec venit, qui praedicaret captivis remissionem.

Deinde contendit reverendus pater, merita augeri relictis indulgentiis. Hoc ego non video, quoniam, ut dixi, relinquentes indulgentias aliquando istam pecuniam in bonis conviviis volunt consumere, non merita augere. Fateor, quod aliquando posset quis facere maius

opus meritorium cum illa pecunia, quam si redimeret indulgentias, casus non oportet specificari.

De iubilaeis celebratis non vult admittere, accessisse consensum fidelium, et suspecta Bonifacii VIII. vita posset suspectum facere eius iubilaeum. Dico, satis magnum esse populi consensum, quod ex toto Christiano orbe in iubilaeo Christi fideles confluunt; quomodo alias consensum fidelium probare possit, non video.

Quod consensu principum dixi pro passagio et cruciatis a summo pontifice indulgentias datas, respondet parum conficere, si hic erraverint in re nihili, quasi nihil esset, Christi fideles decipere et unanimem Christianorum principum et regum consensum in tam fallacem descendere fidelium deceptionem.

Praeterea, nullum, ait, esse decretum expressum et clarum super indulgentiis. Oppono ei declarationem Sixtinam et iam novissime declarationem moderni pontificis, ubi inquit, Romanam ecclesiam, quam reliquae tamquam matrem sequi tenentur, tradidisse Romanum pontificem, Petri clavigeri successorem et Jesu Christi in terris vicarium, potestate clavium, quarum est aperire tollendo illius in Christi fidelibus impedimenta, culpam scilicet et \marginpar{\small[353]} poenam pro peccatis actualibus debitam. Culpam quidem mediante sacramento poenitentiae, poenam vero temporalem pro peccatis actualibus secundum divinam iustitiam debitam, mediante ecclesiastica indulgentia: posse pro rationabilibus causis concedi eisdem Christi fidelibus, qui charitate iungente sunt membra Christi, sive in hac vita, sive in purgatorio, indulgentias ex superabundantia meritorum Christi et sanctorum ac tam pro vivis quam defunctis apostolica auctoritate indulgentiam concedendo, thesaurum meritorum Christi et sanctorum dispensare. Ubi clare determinat cum sequentibus, indulgentias non esse remissiones bonorum operum, sed remissiones poenarum, quae fiant compensatione facta meritorum Christi.

Et quod ratiocinabar, ecclesiam tot annis totam non errasse, respondit, non totam ecclesiam errasse, quia semper interea fuerint, quibus hoc negotium fuerit suspectum. Hinc vulgatum illud: Indulgentiae sunt piae deceptiones fidelium. Ad quod dico: Quando alicuius rei habenda est ratio, non quod quisque dicat, sed quod boni, docti, honesti et optimi statuunt, audiendum est. Tot enim fuerunt

probatissimi et honestissimi viri per totum Christianum orbem in trecentis annis et horum aliqui in catalogum sanctorum relati, non tamen indulgentias hoc pacto improbaverunt, etiamsi abusus eis sicut in omni re displicuerit. Thomas sanctus est, Bonaventura, Albertus Carmelita, Bernhardinus et alii magni viri, Ales, Gerson, Capistranus etc., qui tamen indulgentiarum concessionem non reprobarunt. Quare nihil refert, quod Johanni Wikleff et Ulrico Kalteisen in Anglia, domino Johanni de Wesalia, qui tamen hoc retractavit Maguntiae, et Bohemis indulgentiarum concessiones non placuerunt.

Quare concludamus, indulgentias rationabiliter datas ad dei honorem et fidei exaltationem non esse inutiles Christi fidelibus devotis, non praesumentibus, sed eas in humilitate et gratiarum actione accipientibus, abusibus resectis et repulsis, ut sic potius de tanto munere deo gratias agamus, quam ut id superbe contemnamus, semper in animo habentes, optimam indulgentiam esse veram contritionem. Tamen omnino, si hodie exposita sententia fuit mens reverendi patris, maluissem, ut illud clare dixisset et simplicibus non dedisset occasionem aliud cogitandi, quia nemo est omnium auditorum, qui sensum hodie datum intellexerit in conclusione undecima: Dicere, indulgentias 'esse bonum Christiano, est insanire, sunt enim verissime operis boni vitium. At sum bene contentus de sua expositione, quantum ad me attinet.

\paragraph*{D. D. Martinus.}

Contra solutiones meas opponit egregius D. D. primo, indulgentias non esse remissiones bonorum operum, nec se vidisse, qui hoc assereret. Sequi \marginpar{\small[354]} etiam ad hoc, indulgentias esse vere periculosas, noxias et damnabiles, esse potius eas uno consensu dantium et 'admittentium remissiones poenarum pro peccatis debitarum. Respondeo: Satis admiror, quid haec verba D. doctoris velint, cum sine dubio per poenas pro peccatis debitas intelligant omnes labores satisfactionis, quae tertia pars est poenitentiae. Inter quos certe sunt bona opera, eleemosynae, orationes, ieiunia, vigiliae et similia. Ad quod est expressus textus lib. 5. cap.: Cum ex eo, ubi papa dicit, per indiscretas indulgentias enervari poenitentialem satisfactionem, et omnes bullae indulgentiarum ferme hanc clausulam habent: De iniunctis poenitentiis relaxamus.

Sed age, dato, quod essent poenae et non opera, dico peius esse, remitti poenas quam opera, cum vita passiva et poenosa plus perficiat ac prosit quam mere activa, ut apostolus 2. Cor. 12. [v. 10]: Dum infirmor, tunc fortior sum. Purius est enim opus, quod solo deo operante, nobis patientibus, perficitur, quam quod nobis cooperantibus efficitur, iuxta illud Deut. 32. [v. 39]: Ego percutiam et sanabo, occidam et vivificabo. Quare ad sensum egregii D. D. plura absurda sequuntur ex indulgentiis quam ad meum sensum, et sane nescio, an verissime dicam.

Secundo, non potest concordare, quod dixi, indulgentias solum pigris esse datas et canones similiter tantum pigris impositos. Dico: Non memini me dixisse, canones esse pigris impositos, si autem dixi, sensus meus est, quod strenui et ferventes peccatores plus faciunt, quam canones exigunt.

Tertio etsi fateatur, non esse necessarias indulgentias ad salutem, expedit tamen scire, an promoveant salutem necne. Respondeo: nec hoc est necessarium scire, sed si nocerent saluti; non enim est necessarium scire non necessaria ad salutem.

Item, quid interest, an papa per merita Christi possit remittere? Nescio an intersit. De meritis Christi postea.

Post haec dicit: Meam sententiam destruere dicta conciliorum et indulgentiarum largitionem, quae communi forma concedunt contritis et confessis. Respondeo: Non ideo destruit mea sententia dicta conciliorum, quia sunt contriti et confessi, qui pigri sunt ad solvendas satisfactiones, quos scriptura etiam infirmos vocat, et propter hos damnandae non sunt indulgentiae.

Reprehendit etiam me, quod usurpavi dictum Senecae: Non quis sed quid dicatur, attende, quod in his, quae sunt fidei, attendere oporteat etiam, quis dicat. Respondeo: Hoc non est ad propositum, quia materia indulgentiarum non est materia fidei. Deinde, in materia fidei multo maxime opus est attendere, non quis, sed quid dicatur,

quia non est respectus personarum habendus, sicut manifeste praemonuit Christus Matth. 7. [v. 15] et 24. [v. 5]: Attendite a falsis prophetis, qui veniunt ad vos etc. Item: multi venient in nomine meo dicentes: Ego sum \marginpar{\small[355]} Christus; et Joh. 10. [v. 27]: Oves meae vocem meam audiunt. Proinde etiam summo pontifici non temere credendum est, ut stet regula Johannis apostoli [I. 4 v. 1]: Carissimi, probate spiritus, utrum ex deo sint.

Dicit deinde, periculum esse in errore indulgentiarum ex eo, quod periculosus error habitus sit eorum, qui ordinibus mendicantium ab initio detraxerunt. Quidquid sit de detractoribus, per hoc non probatur error periculosus in indulgentiis, nec mox error est, quod ordines mendicantium vel sacerdotes ecclesiarum errorem iudicaverint. 'Per idem', ad id, quod dicit, concilium non errare, quia regitur a spiritu sancto, addo ego, scilicet in his, in quibus regitur a spiritu sancto, id est, in rebus fidei. Et ut aliquando dicam, non dixi errasse concilium in dandis indulgentiis, sed posse errare.

Transeo illud, quod remissionem peccatorum intellexit remissionem poenarum, cum saepe peccatum pro poena accipiatur, 2. Maccab. 12 [v. 46]. Ego non saepe invenio peccatum pro poena accipi, nisi forte, ubi de Christo scribitur, quod peccata nostra ipse portavit, quod tamen et ipsum non ausim dicere simpliciter pro poena accipi. Sed haec extranea.

Confutaturus illud, quod venias dixi vitium operis, dicit esse subordinationem finium in operibus, ideo si veniae referantur in deum, non vitiant opus. Respondeo: Da exemplum demonstrandi. Dixi enim hodie raros inveniri, qui tantum darent gratis et propter deum, quantum dant pro indulgentiis; ex fructibus ergo cognoscetis eos, imo si audiant melius posse se facere quam redimere indulgentias, probabile est, nullos esse redempturos indulgentias.

De contemptu indulgentiarum dicit: Si fuissem usus hoc moderamine, standum esset meae declarationi, verum vulgaris sermo aliam indicat laicis intelligentiam. Respondeo: Videat diligentius eundem sermonem egregius D. D. et inveniet, me expressis verbis dixisse, non esse contemnendas indulgentias, sed relinquendas liberas, neminem enim absterrendum, sicut neminem urgendum, esse videlicet pigrorum

solatium. Quare sine mea culpa hanc declarationem et satis moderatam sententiam non intelligunt.

Transeo illud de Gersone, ubi dicit: Aliud esse satisfacere in peccatis et redimere indulgentias in peccatis: quia in his obstat forma concessionis, quae solum confessis et contritis permittit indulgentias, illic vero et ipse Johannes Gerson secum sentiat. Dico: Relinquo hoc tamquam extra capacitatem intellectus mei, vere enim non capio, quomodo peccator possit satisfacere in peccatis et non etiam possit recipere licentiam non satisfaciendi, cum plus videatur requiri ad operum perfectionem quam ad omissionem.

\marginpar{\small[356]} De Gregorio, etsi historia non habeatur recepta, tamen famam dicit non permittere negari. Dico: Famae temere nemo credat, et transeo.

Tandem ad id redeo, quod indulgentiae sint satisfactoriae et non remissiones bonorum operum, sed poenarum, fortassis contendimus in aequivoco, quod satisfactorium D. D. appellat, quia permittitur virtute indulgentiarum, ut non satisfaciat, et haec permissio non satisfactionis pro satisfactorio habeatur. Ego satisfactorium appello, quod implet satisfactionem, non omittit.

Adduxit etiam auctoritatem Esai. 61 [v. 1], ubi Christus Lucae 4 [v. 18] legens ex Esaia, dicit: Spiritus Domini super me, eo quod unxerit me, praedicare captivis indulgentiam, ubi Christus remissionem dixit. Addidit Theophilum, qui exponit, quod animae tenebantur etc., donec venit, qui praedicaret captivis redemptionem. Quod illa auctoritas nihil faciat ad rem, ipsa verba et consequentia textus declarant, quia Christus ibidem dicit: Hodie impleta est scriptura haec in auribus vestris; deinde, quia loquitur de vero anno iubilaei, hoc est, de plenitudine temporis, de corona anni benignitatis, et non de remissionibus poenarum, sed de remissione peccatorum, de qua tota scriptura loquitur, cum iubilaeus Bonifacianus tanto post tempore coeperit. Item quod neget augeri merita relictis veniis, transeo, quia non per hoc confutatur, quod aliquando pecuniam eandem volunt consumere in conviviis.

De iubilaeis celebratis dicit, satis magnum esse populi consensum, quod ad indulgentias tot fideles in toto orbe confluunt, excepta

tamen Italia et Roma, nec videt, quomodo alias consensum possim probare fidelium Respondeo: Eum consensum ego probarem esse fidelium, qui re patefacta et veritate indulgentiarum exposita fideles ad confluendum moveret. Nunc vero confluunt ignorantes, quid et quare, cum maior pars credat se rem necessariam et meritoriam facere.

Post hoc de passagiis et cruciatis arbitratur egregius D. D., non esse rem parvi momenti, tot fideles decipi, praesertim magnates. Respondeo: Nec est res magni momenti, quandoquidem et divus Hieronymus Matth. 23. [v. 35] de sanguine Zachariae prophetae dicit: Non damnamus errorem, qui de pietate fidei descendit. Ita et hic error est sine damno animae, citra etiam lucrum pecuniae, quod deferunt pro honore s. ecclesiae rectoribus ecclesiarum.

Quando ego dixi, nullum esse decretum super indulgentiis, opposuit Sixtinam et Leoninam novissimam, ubi declarantur virtutes indulgentiarum. Dico: haec alia quaestio est. Ego dixi, nullum esse decretum, quod praeciperet indulgentias redimere, ideo non esse necessarias.

In fine, ad proverbium illud: Indulgentiae sunt piae fidelium deceptiones, per quod volui, non totam ecclesiam errasse, dicit non esse audiendum, quod quisque dicat, sed quod boni et honesti viri statuant. Nunc vero aliqui etiam [:357] in sanctorum catalogum relati non improbant indulgentias. Respondeo: Nec ego sane. Hoc tamen addo, quod ecclesia accipiens opiniones non facit ideo ex opinionibus veritates. Ideo ad illam Leonis X. declarationem dico, sicut in actis dixi, non probari nudis verbis praesertim hominum, indulgentias esse de thesauro Christi sumptas et solvere poenas iustitiae divinae requisitas, licet opinionem non damnem. Nam, ut dicam, quod sentio, cum Christi merita, sive accipiantur ut suffragia, sive quocunque modo et nomine, nihilominus sunt merita Christi, nec propter usum vel applicationem metamorphosim subeunt. Si autem sunt merita Christi, gratia et veritas sunt, iuxta illud Psal. 25 [v. 10]: Universae viae Domini misericordia et veritas, et Joh. 1 [v. 14ff.]: Gratia et veritas per Jesum Christum facta est; ibidem: Vidimus plenum gratiae et veritatis. Quare etiamsi angelus de caelo aliud persuaserit, non credam esse in manu ullius hominis gratiam et veritatem dare, hoc est, merita Christi dispensare, non tamen resisto pontifici. Sed ad declarationem in actis me refero.

\paragraph*{D. D. Eccius.}

Quia, ut dixi, magna pars dissensionis iacet in eo, an indulgentiae sint remissiones bonorum operum vel poenarum, ubi in meam partem dixi descendere ecclesiae doctores et Romanos pontifices, qui hactenus de indulgentiis scripserunt. Sed hoc repellit acute et docte reverendus pater: Quia remittere poenam 'esset remittere labores satisfactionis', inter quos sunt bona opera, contra decretum Innocentii, can.: Cum ex eo, de poenitentiis et remis., 'ne poenitentialis satisfactio enervetur'. Dico bene remitti labores satisfactionis, sed non propterea bona opera, quamvis enim oratio, ieiunium etc. sint satisfactiones, tamen sunt bona opera, etiamsi his non satisfaciamus. Accedit, quod satisfactio non solum fit labore bonorum operum, sed etiam satispassionis. Quare plerique volunt, non debere cogi filium confessionis ad satisfactionem in se recipiendam; iuxta illud commune: Melius est, ut confessor filium cum parva poenitentia mittat ad purgatorium, quam cum magna ad infernum. Nec mirum, si pontifex plenarias resecuit indulgentias ad satisfactiones poenitentiales eludendas, si passim a praelatis fierent. Sic particula bullae de iniunctis poenitentiis, id est, quae secundum iustitiam divinam iniungi deberent, ut doctores ecclesiastici de indulgentiis scribentes declarant.

Porro, quod argumentatur, peius esse remittere poenas quam opera, non accipio, quia pro tanto poenae prosunt, ut homo instruatur bene operari, et sic virtus in infirmitate perficitur, sed hoc ad iudices. Et quia errare in indulgentiis esset animabus 'periculosum', ut hodie induxi, ideo ecclesia non errat, si solum pigros contritos vult recipere indulgentias, certe ad Carthusianos venire non deberent.

\marginpar{\small[358]} In materia fidei dicit maxime attendendum, quid dicatur. Ego autem auctoritatem reputo in theologia potentissimam, et quod audire praecipimur sedentes super cathedram Moysi etc.. Ideo quis dicat, attendere debemus, ut captivemus intellectum, iuxta praeceptum apostoli, in obsequium fidei, propterea concilia habita sunt, ut eorum auctoritate errores exterminarentur.

De mendicantibus transeo.

Quod vero ait, concilium in his, in quibus regatur a spiritu

sancto, non errare, dico ego, antiquorum patrum secutus sententiam, semper debere praesumi concilium regi a spiritu sancto, quamdiu non constat de opposito, scilicet de illegitimitate congregationis.

Sed adhuc fortiter stat reverendus pater in vitio boni operis, quod dans pecuniam propter indulgentias alias non daturus, quasi hoc sit vitium boni operis. Quaererem a 'reverendo patre', si ipse 'cum suis fratribus' habeat anniversarium alicuius 'principis' vel alicuius nobilis cum triginta missis, quia dantur bonae praesentiae, alioquin non habiturus anniversarium: an hoc sit vitium boni operis? [Martinus dubitat. Eccius: Et sic non erit salva omnis caro tam in cuculla quam extra.]

De indulgentiis ab eo non contemptis iudicent, qui sermonem vulgarem legerint.

Porro reverendus pater factus interpres meus dicit, indulgentias forte sic esse satisfactorias, quia propter ipsas aliquis 'satisfacere non teneatur'. Non iste est sensus meus, sed quia homo deberet pro peccatis etiam contritis satisfacere de suo et per indulgentias satisfacit de alieno, quia summus pontifex dat ei de thesauro ecclesiae, unde solvat, ut sic, iuxta communem patrum sententiam, etiam in indulgentiis peccatum non maneat impunitum, salvo eo, quod etiam de proprio satisfacere non posset, non accedentibus Christi meritis.

De consensu fidelium credo neminem esse tam stultum, qui credat, rem esse necessariam indulgentias in iubilaeo, et si existimet iter Romanum sibi esse meritorium, nec hic errat, stat enim opus, per quod quis consequitur indulgentias, esse meritorium et tamen ipsas indulgentias non esse meritorias.

Porro, de verbis Sixtinis et moderni pontificis non facio vim, quia cathedram Petri sequor et eius sententiam, qui in ea sedet, quamdiu in haeresim, quod absit, prolapsus non fuerit. Scio eum non probare, sed definire, quem ut Christi vicarium definiendo arbitror esse indefectibilis fidei. Sed haec omnia eorum committo iudicio, quorum interest, paratus errata expungere, si quae ostendentur.

Addidit D. D. Martinus: Et ego.

\marginpar{\small[359]} Crastina hora septima ante meridiem inter eosdem E. D. doctores continuabitur disputatio tractabiturque materia de poenitentia.

Die Martis duodecima mensis Julii, hora septima antemeridiana.

\paragraph*{D. D. Eccius.}

Contra conclusionem reverendi patris et partem sermonis sui de poenitentia pro defensione optimorum patrum praedicantium intendo probare, poenitentiam veram incipere a timore etiam poenae, et quod bene quis se praeparet discutiendo, recogitando etc.

Primo, quia talem modum praedicandi observat Dominus Jesus et eius praecursor s. Johannes. Luc. enim 15. filius prodigus gerit typum poenitentis secundum Augustinum de quaest. evang., Ambrosium libro 2. de poenitent. cap. 3., Chrysostomum, Hieronymum et alios, sed eum ita proponit nobis Christus: in se conversus dixit: Quanti mercenarii in domo patris mei abundant panibus, ego autem hic fame pereo, surgam et ibo ad patrem meum et dicam illi: Pater, peccavi in caelum et coram te etc. Hic Dominus Christus modum poenitentis describens, exponit primo motum poenitentem praemiorum magnitudine, scilicet »abundant panibus«, et timore poenae, scilicet »hic fame pereo«, quibus gradibus evectus poenitentiam veram coepit meditari, scilicet »et dicam: Pater, peccavi« etc. Et ut Basilius explanat: Tres sunt ibi poenitentiae gradus, mercedis spes, suppliciorum timor et bonitatis paternae sincera dilectio. Et sic, antequam fiat amicus dei, qui extitit inimicus, efficitur prius servus.

Sic Johannes incepit a timore: Progenies vel genimina viperarum, quis demonstravit vobis fugere etc.?

Id etiam in toto actum est veteri testamento, ut timore poenarum ad mandatorum obedientiam homines traherentur, ut in legalibus, historiis et prophetis hoc ubique reperitur. Sufficiat adduci unum locum, Psal. 89. [v. 31ff.]: Si autem dereliquerint filii eius legem meam et in iudiciis meis non ambulaverint, si iustitias meas profanaverint et mandata mea non custodierint, visitabo in virga iniquitates eorum et in verberibus peccata eorum; cum concordantiis. Quare praedicatores sacram imitati scripturam hactenus bonum habuerunt modum de poenitentia paranda, recogitando gravitatem peccatorum, aeterni supplicii etc.

Accedit ratio Dionysiaca: Quia sic deus res administrat, ut ima ad superiora per media ducat. Quare cum peccator sit in imo, quia non est dignus pane, quo vescitur, volens ascendere ad superiora gratiae hoc faciet per medium timoris.

Bene esset perfectionis, quod aliquis immersus luto \marginpar{\small[360]} peccati posset se mero obtutu dilectionis dei et amore iustitiae erigere ad gratiam capessendam; sed quis est hic et laudabimus eum?

Imo b. Augustinus approbat modum praedicandi nostrae tempestatis, docens expresse, nullum pervenire ad charitatem et ad veram gratiam, nisi praecedente timore, non filiali, sed etiam servili. Aug. testatur tract. 9 super Joh.: Ergo incipiat timor, quia initium sapientiae timor Domini. Cum autem coeperit charitas habitare, pellitur timor, qui ei locum praeparavit. Quantum enim illa crescit, ille decrescit, et quantum illa fit interior, timor pellitur foras; maior charitas, minor timor; minor charitas, maior timor. Si autem nullus timor, non est, qua intret charitas (dat comparationem), sicut videmus per setam introduci linum, quo suitur, seta prius intrat, sed nisi exeat, non succedit linum: sic timor prius occupat mentem, quia ideo intravit, ut introduceret charitatem. Allegat illud psalmistae [30 v. 12]: Convertisti luctum meum in gaudium. Et infra: sententia dicta est de scripturis, nam qui sine timore est, non poterit iustificari. Opus ergo est, ut intret primo timor, per quem veniat charitas. Timor medicamentum, charitas sanitas. Ex istis apertissime liquet, quia poenitentia est medicinalis et timor est medicamentum secundum Augustinum et charitas non intrat nisi per timorem, ideo bene praedicatur, poenitentiam a timore incipere.

Et non accipio, quod reverendus pater ab initio resolutorii refert reverendi patris Staupicii vocem quasi caelitus demissam, poenitentiam incipere ab amore et dilectione. Nam huic expresse reluctatur b. Ambrosius, bonus animarum medicus, in epistola ad Studium: Ubi poena praescribitur, debet esse poenitentia peccatorum; ubi remissio donatur, gratia est; praecedit poenitentia, sequitur gratia. Neque ergo poenitentia sine gratia, neque gratia sine poenitentia. Debet enim poenitentia prius damnare peccatum, ut gratia possit id abolere. Illud

idem testatur Chrysostomus toto libro de cordis compunct. et Homilia 80. de poenit. et Sermone 29.

Addo etiam Isidorum libro 2. de summo bono. ca. 12.: Compunctio cordis est humilitas mentis cum lacrimis, exoriens de recordatione peccati et timore iudicii, et toto libro de contritione cordis et speculo peccatorum Augustinus nihil aliud agit, imo hortatur peccatorem volentem poenitere, ut recogitet tres abyssos, scilicet peccatorum suorum, poenarum et iudiciorum dei. Sic b. Bernhardus super Canticis ser. 16. et aliis pluribus locis, Gregorius in moralibus lib. 2. et 5., et item super ca. 29. Job, libro quoque 2. super Ezech., ubi Homilia 19. inquit: Scriptum est: Initium sapientiae timor Domini. Constat procul dubio, quia a timore ad Dominum ascenditur, non autem a sapientia ad timorem reditur. Propheta ergo de caelestibus ad ima loquebatur, coepit magis a sapientia et descendit ad timorem. Sed nos, qui a terrenis ad caelestia tendimus, eosdem \marginpar{\small[361]} gradus ascendendo numeramus, ut a timore ad sapientiam pervenire valeamus. Idem lib. 1. super Ezechiele dixit, duas esse pennas, scilicet timorem et poenitentiam, quae tegere corpora dicuntur, id est, peccata.

Ex quibus omnibus, quoad unam particulam de timore, volo habere conclusum, bene praedicasse nostrae tempestatis et superioris praedicatores, quod poenitentia a timore incipiat: et quod nituntur sermonibus suis in populo timorem dei seminare, ut sic semen diaboli extirpetur, iuxta illud Origenis lib. 3. super Job: Bonus est timor poenae et iudiciorum, quem nisi daemon expulerit, non poterit seminare semen peccatorum. Ex quo dein consectarium est, poenitentiam non incipere ab amore et dilectione iustitiae, quamvis si sic inciperet, faterer eam esse laudabiliorem et perfectiorem, quam quod a timore poenae incipit. At fragilitas nostra id non patitur, cui Dominus Jesus et praedicatores condescendentes timorem praedicant, tamquam gradum, quo verum iustitiae amorem consequamur.

\paragraph*{D. D. Martinus.}

Non est iste modus scripturas divinas feliciter intelligendi vel interpretandi, si ex diversis locis diversa decerpantur dicta, nulla habita ratione vel consequentiae vel collationis. Imo iste est canon errandi vulgatissimus in sacris literis. Oportet ergo theologum, si nolit errare, universam scripturam ob oculos ponere et contraria contrariis conferre, et sicut

duo Cherubim adversis vultibus utriusque diversitatis consensum in medio propitiatorii invenire, alioquin cuiuslibet Cherubim vultus longe divertet sequacem oculum a propitiatorio, id est vera Christi intelligentia.

Proinde nunquam mihi egregius D. D. remotior visus est a sacris literis, quam hodie, atque eo magis, quod concludendo fatetur, laudabiliorem et perfectiorem esse poenitentiam, si ab amore iustitiae inciperet, sicut ego sapio, quam si a timore poenae inciperet. Quasi non conandum sit, ut laudabiliter et, ut Johannes ait, dignos fructus poenitentiae agamus. Nam hoc prorsus non accipio, quod dicit, fragilitatem nostram obstare, quominus ab amore iustitiae poenitentiam incipiamus. Si secundum fragilitatem nostram agendum est et scriptura exponenda, nunquam poenitebimus, sed de die in diem peiores fiemus. Quare antequam respondeam suis obiectis, primum declaro me ipsum.

Omnis bona vita necesse est ut instituatur per aliquam legem, ideo lex principium est poenitentiae et cuiuslibet boni operis. Quare in poenitente ante omnia oportet vel revelari vel suggeri legem, contra quam fecerit, et secundum quam facere debet. Lege autem manifestata aut in memoriam revocata mox sequitur augmentum peccati, si desit gratia, quia naturaliter odit voluntas legem, ut sunt expressa testimonia Pauli ad Rom. et ad Galat. Rom. 5. [v. 20]: Lex intravit, ut abundaret delictum, Galat. 3. [v. 19]: Lex \marginpar{\small[362]} propter transgressiones posita est. Ideo b. Augustinus libro de spir. et lit. c. 3: Neque enim liberum arbitrium ante gratiam quicquam valet, nisi ad peccandum, non autem ad poenitendum, ut D. D. dicit (hoc enim Pelagianum est). Sequitur Augustinus: Et cum coeperit non latere, quod faciendum est, nisi spiritus sanctus diffuderit charitatem in cordibus nostris, non diligitur, non suscipitur, non bene vivitur. Ibi clarissime dicit Augustinus, quod lex dei non potest diligi, nisi accepta gratia spiritus. Si autem non diligitur lex, contrarium eius peccatum non oditur. Ergo impossibile est poenitere ante dilectionem legis. Hoc est, quod Roma. 4. [v. 15] apostolus vult: Lex iram operatur, hoc est, monstrat peccatum, sed non dat gratiam, ut odiatur peccatum. Ideo manet odium legis et dilectio peccati, quantumlibet per increpationes forinsecas aut intrinsecas homo

concutiatur timore servili. Nam etsi abstinet ab opere peccati, non tamen abstinere potest ab amore peccati. Hoc et Christus Johan. 6. [v. 44] docet, ubi dicit: Nemo venit ad me, nisi pater meus traxerit eum. Concedo ergo, quod lex, recordatio peccatorum, intuitus poenarum possunt terrere peccatorem, sed nunquam facient poenitentem.

Respondeo ergo ad obiectionem primam, de filio prodigo Luc. 15, ubi incepit poenitentiam a magnitudinis praemii recogitatione dicens: Abundant panibus etc. Dico, quod hic prodigus filius incepit vere ab amore iustitiae, quia conversus in se cognovit primo bonum et ex cognito bono intellexit suum malum. Hanc autem conversionem sui in se ipsum non habuit ex fragilitate sua aut ex timore poenae, sicut ipsemet D. D. hic dicit, quod timor poenae post praemiorum magnitudinem moverit. Ergo non incepit a poena et timore. Habuit autem eam ex trahente intus patre et dilectionem inspirante paternae domus, dum dixit: Quanti mercenarii in domo patris mei. Nam cum prius vixisset in peccatis, bonum nec cognovit nec dilexit, nec peccatum odivit: cum tamen non posset ignorare peccatum. Ergo necessaria fuit alia affectio, id est, amor boni.

Secundo induxit Johannem baptistam, Luc. 3 [v. 7], quod inceperit a timore, dicens: Quis monstravit vobis fugere etc.? Respondeo: Aliud est praedicare poenitentiam, aliud incipere poenitentiam; aliud est praedicare bonum opus, aliud incipere bonum opus. Praedicator monet, terret, allicit etc., sed nihil sequitur, nisi gratia moverit voluntatem.

Idem dico ad Psal. 89 [v. 33]: Visitabo in virga iniquitates eorum etc. Percuti potest peccator, sed nisi gratia cooperetur, nihil proficitur, ut Jere. 5 [v. 3]: \marginpar{\small[363]} Percussisti eos, et non doluerunt; et Esai. 1 [v. 6]: A planta pedis usque ad verticem non est in eo sanitas: super quo vos percutiam ultra? quia nihil profecit percutiendo.

Et satis admiror egregium D. D., quod oblitus luminis naturae Aristotelis audeat statuere alicuius virtutis initium in timore urgente, cum ille tot verbis iam receptissimis 3. Ethicorum persuadere conetur: Oportere esse voluntarium opus bonum et libera voluntate fieri. At voluntas certe aut amor est aut soror amoris. Inde etiam dicitur, oportere opus bonum fieri per liberum arbitrium. Addo et illud, quod Christus nunquam peccatores coegit timore ad poenitentiam, sed suaviter allexit quoscumque vocavit, ut Zachaem, Magdalenam, apostolos et omnes. Sicut et apud Jere. cap. 31 [v. 3] dicit: Perpetua charitate dilexi te, ideo attraxi te miserans.

Dico ergo, quod timor Domini quidem necessarius est, sed filialis, quia sine amore impossibile est ferre conversionem sui, in qua terretur, conteritur et humiliatur peccator, iuxta illud I. Reg. 2 [1. Sam. v. 6]: Dominus deducit ad inferos et reducit. Credo autem, etiam ipsius egregii D. D. sententiam et omnium scholasticorum doctorum mecum esse et contra obiecta eius pugnare, cum omnes consentiant, contritionem oportere fieri in charitate, si debet esse bona et meritoria poenitentia. Quod ego sane intelligo contritionem fieri movente et imperante charitate, ut sic sit voluntaria, hilaris, amorosa poenitentia. Ideo licet Johannes increpaverit Judaeos et terruerit, non ideo sequitur, poenitentes incipere etiam a terrore. Aut si incipiunt a terrore, nisi accesserit gratia, hypocritae sunt verius quam poenitentes.

Adiecit D. D. rationem, quod deus res sic administrat, ut ima per media ducat ad superiora, quod aptavit ad timorem, volens peccatum esse imum, timorem medium, amorem superius, quod ego transeo et non accipio.

Dicit etiam, esse perfectionis, si homo possit se mero dei obtutu et amore iustitiae erigere ad gratiam capessendam. Sed quis est hic, et laudabimus eum? Respondeo: Nec timore nec amore potest se erigere homo ad gratiam capessendam, sed gratia praevenit et movet ad merum dei obtutum et amorem iustitiae.

Ad Augustinum, qui doceat, praecedere timorem ante gratiam et intrante charitate pelli timorem, dico: Si recte intelligatur, admitto, hoc est, quod poenitentia nondum est incepta, quando timor praecedit charitatem. Sed intrante charitate incipitur poenitentia, id est, amor iustitiae et odium peccati. Si autem charitas non intraret, timor non operaretur nisi maiora peccata. Quare ista similitudo setae et lini probatur mihi, modo non intelligatur, quod timor inducat charitatem, quod ex ipsius D. D. verbis capio, quod nisi exeat seta, \marginpar{\small[364]} non succedit linum, hoc est, nisi timor, qui impedit poenitentiam veram, pellatur per intrantem charitatem, nunquam vere poenitetur, ut sic dixerim.

Ambrosium inductum in epistola ad Studium, quod praecedit poenitentia et sequitur gratia, et alia, item Isidorum de summo bono etc., item tres abyssos recogitationes peccatorum, poenarum, iudiciorum), Bernhardum super canticis, Gregorium in Moralibus) et alios patres, qui docent a timore) ascendendum ad sapientiam et poenitentiam a timore incipiunt, libenter admitto, sed non contra apostolum Paulum de lege et timore legis docentem intelligo. Dico et ego, habita charitate simul moveri hominem ad timorem dei, et sic incipi) poenitentiam a timore in charitate, alioquin stat firma sententia, quod timor poenam habet et) bonum non operator, sed odit legem.

Non ergo egregius D. D. concludit, poenitentiam non incipere ab amore iustitiae, quantumlibet bonus sit timor poenae ex origine). Non enim timore poenae, sed timore dei poenitendum est, quod ille sit servus, non mansurus in domo, hic autem filius et heres. Ideo et illud Proverb. 1 [v. 7]: Principium sapientiae timor Domini, non admitto intellectum de timore poenae, qui ante gratiam torquet hominem infructuose, cum expresse dicat: Timor Domini, non timor poenae. Timor poenae potius est principium insipientiae. Videat ergo egreg. D. D., ut non in unum chaos confundat timorem servilem et filialem), ne sibi ipsi) scripturae et patrum intelligentiam praecludat.

\paragraph*{D. D. Eccius.}

Reverendus pater fuco verborum conatur se evolvere ex tam expressis sacrae scripturae et s. patrum auctoritatibus. Et ut hoc vobis persuaderet, ausus est dicere, me longe fuisse alienum a sacrae scripturae intelligentia, hoc loco diverticulum quaerens de Cherubim se respicientibus. Sed iudicent, quorum interest, uter rectius) sentiat de sacra scriptura. Tamen ut suas solutiones impertinentissimas excludam, duo frequenter adducit, quorum in sermone recitando viam communem non meminit), nec ullus praedicatorum, quod meminerim), negavit, et hunc quoque timorem, in quantum disponit ad veram poenitentiam, praeveniri inspiratione divina. Nam indubitatum est apud Christianos contra Pelagii) perfidiam, quod salutis nostrae initium) deo inspirante habemus. Quare non fuit necessarium, hoc afferre vel ob hoc modum praedicandi reprehendere.

In uno tamen videtur mihi aequivocatione falli, quod hanc gratiam, qua deus praevenit corda hominum movendo, credit esse charitatem, cum sit aliud gratuitum dei munus. Et per hoc resoluta est b. Augustini de spir. et lit. sententia contra improbos Pelagianos.

Alterum, quod dicit, timorem esse infructuosum, nisi accedente charitate, quis hoc unquam vel scholasticorum vel praedicatorum negavit, qui omnes apostoli Pauli ad Corinth. de charitate sequuntur sententiam?

\marginpar{\small[365]} Porro, quod laudabiliorem dixi poenitentiam, quae ab amore incipiat, assumpsit nos debere iam) facere iuxta illud Johannis: Facite dignos fructus poenitentiae [Matth. 3 v. 8]. Nam si secundum fragilitatem nostram esset agendum, nunquam poenitebimus. Dico nos quoque posse facere dignos fructus poenitentiae, etiam si a timore inceperimus et ad charitatem perveniamus.

Et miror, quod reverendus) D. pater vult nos facere angelos et fragilitatis nostrae oblitos), cum b. Gregorius lib. 2. super Ezechielem expresse declaraverit), prophetam descendisse de sapientia ad timorem, sed nos eosdem gradus ascendendo a timore ad charitatem pervenimus, sicut multa alia fragilitati nostra condonantur.

Mentem suam declaraturus ostendit, quomodo dilectio legis praecedere debeat poenitentiam: quia per legem abundavit delictum [Röm. 5 v. 20], et ad Gal. 3. [v. 19] etc.). Dico: Verum est, ut s. Augustinus contra Faustum lib. 22. testatur, peccatum factum vel dictum) contra legem ideo non esse) peccatum, si non esset lex prohibens, sed sicut legis transgressione abundavit) delictum, ita legis observantia augetur meritum. Quare omnino ista diverticula ad propositum non conferunt: et solutiones datas esse contra sacrae scripturae intelligentiam), post prandium manifestabimus.

Hora secunda pomeridiana continuabitur disputatio.

Hora igitur secunda eiusdem diei

D. D. Eccius prosecutus est replicam suam hodie coeptam).

Reverendus pater ad illud de filio prodigo) respondit, praecessisse amorem iustitiae, iuxta illud: »Et conversus in se«, quia memoriam paternae domus habuerit. Sed hoc non diluit obiectum, quia conversio

illa fuit facta contemplatione poenae, dum neminem haberet, qui eum siliquis satiaret, durante enim pecunia non convertebatur in se, sed famelicus conversus in se etc.). Accedit, quod b. Augustinus ait, verba eius fuisse verba poenitentiam meditantis, nondum agentis. Quodsi conversus in se fuisset amore iustitiae, iam poenitentiam incepisset contra Augustinum. Quare solutio sua propositum non infringit et Basilio manifeste reluctatur.

Ad Johannem respondit, aliud esse incipere, aliud praedicare, et alias dixit, praedicatores terrere, sed non facere poenitentem per poenarum cruciatum). Iterum evanida est solutio, quia si Johannes praedicavit, certe intendebat fructum suae praedicationis, et quia insinuabat) eis timorem, signum est hoc pacto poenitentiam auspicari. Et nostri praedicatores praedicant, quomodo incipiant, scit deus.

Porro quod poenitentiae initium) gratiae tribuit, et hi) praedicatores et doctores nunquam negaverunt, quia deus aspirando praeveniat).

Praeterea Aristotelem non vult acceptare in scholis theologicis, et tamen audet eum mihi opponere. Sed dico, constantem esse ethnicorum sententiam, \marginpar{\small[366]} sistendo in timore poenae non esse perfectam virtutem), iuxta illud: Oderunt peccare mali formidine poenae, oderunt peccare boni virtutis amore). Nec opus virtutis debet esse coactum, sed liberum. Aliud tamen est liberum arbitrium induci, aliud) liberum arbitrium cogi. Fateor, Christus benevolis verbis vocavit apostolos, Zachaeum, Magdalenam. Sed ipsum aliquando durius vocare, factum Pauli nobis innuit, de quo Augustinus ait 23. quaest. 4. can. Quis nos): Cui Christus vim intulit? quem coegit? et in evangelio de vocatione ad coenam dicit: Compelle intrare. Sicut pluribus verbis b. Gregorius in homilia prosequitur. Taceat ergo, quaeso, reverendus pater, solum benevolentem afferens) vocationem.

Praeterea solum timorem filialem necessarium arbitratur poenitenti, de quo plurimum miror, cum hodie me ad veram) scripturarum intelligentiam monuerit, quod se ipsum primo non emendavit, nam et b. Augustinus in loco hodie allegato super Johannem loquitur de timore servili, quem foras) mittit charitas, et de illo dicit loqui Sapientem: Initium sapientiae timor Domini. [Prov. 1 v. 7] Et ita de timore servili intelligit) glossa super Psal. [111 v. 10]: Initium sapientiae etc.) Metus iudicii ianua est conversionis ad deum, et est glossa Cassiodori. Sic glossa super illud apostoli ad Roma. 8 [v. 15]: Non enim accepistis

spiritum servitutis in timore, dicit timorem servilem esse bonum et a deo. Sic Paulus testatur illum bonum 1. ad Timoth. 5 [v. 20]): Peccantem coram omnibus argue, ut et ceteri timorem habeant. Ita Augustinus super Psalmis ait, et Longobardus eum renarrat) in 3.: Timor servilis est, cum per timorem gehennae) continet se homo a peccato, quo praesentiam iudicis et poenas metuit. Et infra: Bonus est iste timor et utilis licet insufficiens, per quem paulatim fit consuetudo iustitiae. Beda quoque exponit: Initium sapientiae timor Domini, scilicet servilis.

Quare aliud quaerat sibi patrocinium et aliam det solutionem, quando ista) s. patribus non concordat. Quodsi nolit dictis sanctorum stare, habebit adversum se sacrae scripturae locos. Nam dum Domini timor) est initium sapientiae, adhuc charitas foras) pellit timorem, sed timor filialis iuxta David in aeternum permanet et in saeculum saeculi. Quare alius est timor, quo sapientia initiatur, alius est, qui cum aucta permanet. Unde cum reverendus pater negat timorem inducere charitatem, expresse s. Augustino contradicit, qui in loco allegato hoc totum ait), timorem scilicet servilem esse initium sapientiae et per charitatem iterum expelli. Subscribit sententiae Augustinianae mellifluus Bernhardus Ser. 52. super Canticis, quem brevitatis gratia non adduco. Sed etiam plus dico, charitatem non sic expellere timorem, sicut \marginpar{\small[367]} unum incompatibile) aliud, sed successu ad augmentationem gratiae minui) timorem servilem et aucta gratia timorem omnino auferri.

Nec refert, quod ponderabat hodie D. pater: Timor Domini, non timor poenae est initium sapientiae, quoniam et timor poenae et timor reverentiae uterque est Domini et uterque est a Domino. Sola media timendi sunt diversa, ut Beda et alii per Longobardum inducti testantur praeter Augustinum prius citatum. Et satis liquet ex verbis Christi, si nihil aliud haberemus, Matth. 10 [v. 28]: Nolite timere eos, qui occidunt corpus, animam autem non possunt occidere, sed potius timete eum, qui potest animam et corpus perdere in gehennam). Si enim timor servilis esset damnabilis, sicut dicit reverendus pater, cur Christus nos ad eum invitaret? et Augustinum et alios ita intellexisse, confero in iudicium eorum, quorum intererit.

Insuper dicit, scholasticos omnes concordare, quod contritio nullius sit utilitatis, nisi fiat in charitate. Hoc quidem verum, sed secundum Augustinum ad charitatem non pervenitur, nisi per timorem. Timor medicamentum est, charitas sanitas. Unde non accedente charitate omnes

bene noverunt timorem illum insufficientem, propterea in via illum constituunt), non in termino.

At) rationem ex b. Dionysio sumptam transivit, sed non aperuit D. pater, cum peccatum sit imum et charitas superius, quod sit medium perveniendi) a peccato in charitatem, et profecto aliud dare nequit), quam Basilius, Beda, Augustinus, Bernhardus et ceteri), timorem.

Deinde) assumit gratiam praevenire timorem et amorem, quod si loquitur de gratia motionis divinae, qua deus nos aspirando praevenit, fateor quidem); si autem loquitur de gratia charitatis, illud non accepto, quia illius), quae recte) sapientia in scripturis appellatur, initium est timor Domini, quamvis timor iste absque charitate sit infructuosus.

Porro, quod Ambrosium, Gregorium et alios doctores admittit, sed non contra apostolum de timore legis, nescio quas offundit) nebulas. Dicat clare, an) isti s. patres per me allegati repugnent dictis apostoli vel minus. Si repugnant) dicto apostoli, hoc ostendat, quod nunquam fieri posse arbitror, cum in sacris scripturis fuerint exercitatissimi et spiritu sancto repleti, apostolum Paulum aeque bene intelligentes sicut nos. Si non repugnant dictis Pauli, stet eorum sententiae et praedicatores modum contritionis et poenitentiae praefatum docentes non improbet, nec ab hominibus peccatoribus huiusmodi timorem servilem utilem et quasi necessarium medium excludat.

Non relevat eum, quod existimat habita charitate moveri mentem ad timorem, quoniam hoc esset cancrino more retrogradi) contra b. Gregorii sententiam lib. 2. homilia 19. super Ezechielem). Constat procul dubio, quia a timore ad sapientiam ascenditur, non autem a sapientia ad timorem reditur. Habet ergo timor poenam, quae charitate aucta minuitur et totaliter absorbetur.

\marginpar{\small[368]} Praetereo ob temporis angustiam alia in eodem sermone contenta, quae tamen omnino venirent discutienda, nisi hodiernus dies praefixus esset ad hanc materiam terminandam, sicut est punctus de sigillata confessione peccatorum in specie), uti meminit Chrysostomus super caput 12. Matth. de puncto, quod scholastici ponant tres partes poenitentiae, scilicet contritionem, confessionem et satisfactionem, quas prius posuit Chrysostomus sermone 29. de poenitentia, de duplici modo confitendi sacerdoti magna et alia deo cum aliis punctis, quae omnia

prolixe possent) tractari et excuti. Sed iacto iam fundamento totum simul sermonem et ea, quae adduxi pro parte mea, confero in iudicium iudicum eligendorum.

\paragraph*{D. D. Martinus.}

Sperabam egregium D. D. confutaturum meas solutiones, praesertim ea, quae pro fundamentis ex Paulo adduxeram) de lege, quae ante charitatem non operatur nisi iram et auget peccatum, taceo, quod disponat ad gratiam, quemadmodum et) timor servilis fructus legis iram operatur et auget peccatum. At ipse haec fortiter transiliens, eandem cantilenam recantavit, ac nobis haec non fucum), sed medullam scripturae contendit persuadere, discurram per singula.

Primum dicit), non fuisse necessarium afferre, quod initium salutis nobis) sit ex deo inspirante, ubi et timorem disponentem dari intelligit. Respondeo: Haec plane omnia sunt erronea propter verbum Pauli, qui dicit: Impossibile esse legem impleri, imo peccata non augeri, nisi spiritus sanctus diffundat charitatem in cordibus nostris. [Röm. 8 v. 3 ff.] Has auctoritates tam expressas oportuit confutatas, et illam b. Augustini cap. 3. de spir. et lit.: Dum) coeperit non latere, quomodo vivendum est, nisi gratia donetur, non suscipitur, non bene vivitur. Vadat ergo cum suo servili timore, qui non operatur nisi odium legis et dei et cum iniuria vocatur disponens ad gratiam.

Invenit etiam aequivocationem gratiae, aliam esse charitatem, aliam donum, quo primo) movemur. Transeat illa) distinctio, nihil ad propositum, est elusio) verborum Pauli. Clarissimus est textus Pauli, nisi gratia et charitas faciat nos diligere legem, lex semper iram operatur. Hoc donum autem, quo primum movemur, non diliget) legem, sed charitas spiritus.

Tertio concessit, timorem esse infructuosum sine charitate, et hoc nullum negasse. Respondeo: Cur ergo docent infructuosas poenitentias et resistunt mihi fructuosam docenti? Relinquo, confutationem hanc esse omnino subversam), cum tot opusculis Augustinus), et si Augustus non faceret tot fulminibus), tamen unus Paulus concludit, opera quaecunque ante charitatem esse peccata et damnabilia \marginpar{\small[369]} et

indisponentia ad gratiam. Et per haec dicta ferme ad omnia sua replicata responderi potest. Ipse quidem pro more suo dicta scripturae divinae temperat ad dicta patrum, imo trahit ad intelligentiam suam, quam habet in patribus, cum potius contra scripta patrum debeant conferri et iudicari ad dicta scripturae.

Quarto non accipit, quod assumpsi, nos debere facere laudabiliorem illam poenitentiam, sed dicit, nos posse etiam dignos fructus poenitentiae facere, si a timore inceperimus. Hoc secundum Augustinum) sic expono: Si peccata peccatis addiderimus, cum, ut saepius iam dictum est, quidquid ante sanatam) per gratiam voluntatem fit, fructus malae arboris sit, quae non potest fructus bonos) facere. Ergo per timorem servilem nunquam pervenitur ad charitatem.

Quinto, innixus Gregorio super Ezechielem, qui descendat de sapientia ad timorem, sed nos contra ascendere) a timore ad charitatem docenti, sic respondeo, quod divus Gregorius non debet excludere charitatem nec exclusit a timore incipiente poenitentiam, sicut) D. D. intelligit; sicut nec a charitate excluditur omnis timor servilis, praesertim in hac vita, cum charitatis officium sit expellere timorem servilem, atque id tota vita, et inducere timorem filialem.

Sexto, ad auctoritates Pauli, quod lex propter transgressiones posita est et auget delictum, dicit, nihil esse ad propositum et me diverticula quaesivisse, adducens Augustinum contra Faustum lib. 22: Peccatum esse dictum vel factum contra legem dei etc.); nescio quid velit D. D. Non est hic disputatio, quid sit peccatum, ideo superfluis verbis hoc tempus perditum est. Scopus iste versatur, quod auctoritate Pauli timor servilis non possit esse nisi peccatum et augere peccatum per virtutem) legis, si desit gratia. Et sic auctoritates Pauli non sunt diverticula, sed fulmina conterentia sententiam D. doctoris usque ad pulverem. Non ergo ego ex hominibus angelos facio, oblitus fragilitatis, sed doceo, ne deos faciamus ex peccatoribus, dum obliviscuntur fragilitatis suae, qua non possunt nisi malum facere ante gratiam.

Quod conversionem filii prodigi factam dicit contemplatione poenae, nego. Ad probationem, quod neminem haberet, qui eum siliquis satiaret, dico: Nisi intus fuisset tractus et in se conversus, potius fame mortuus esset, quam rediisset), ut stet sententia: Nemo venit ad me, nisi pater meus traxerit eum [Joh. 6. v. 44]).

Quod Augustinus dicit), verba eius fuisse verba poenitentiam meditantis, nondum agentis, spero, quod pro me faciat, quando meditare

poenitentiam, praesertim medullitus in corde, sit certe incipere poenitentiam. Ergo illud »nondum agentis« oportet de perfectione ad extra intelligi.

Idem ad Basilium dicetur).

De verbo Johannis baptistae dicit, evanidam esse solutionem, quod dixerim, aliud esse poenitentiam docere, aliud incipere. Respondeo: Salva \marginpar{\small[370]} reverentia egregii D. D.), non videtur Paulum intelligere nec virtutem) legis cognovisse. Lex enim docet quidem sancta), iusta et bona, sed sola gratia incipit, facit et perficit ea. Ergo etiam si) Johannes timorem docuisset esse initium poenitentiae, non ideo sequitur, poenitentiam incipere a timore. Sicut si insinuem quodlibet bonum opus per terrores et minas, non ideo incipit bonum opus) a terrore et minis, sed a charitate.

Indignatur, quod Aristotelem ei obiecerim in scholis theologicis, quem tamen non acceptarim. Fateor, minor est, quam ut valeat in theologica schola, cum seductor fuerit scholasticorum doctorum. Sed satisfacere volui conclusioni meae, quod contra sacrum suum Aristotelem desipiunt, qui poenitentiam a timore et non libera voluntate incipiunt. Nam distinctionem) de libero arbitrio coacto et 'inducto' transeo. Non enim verum est, quod liberum arbitrium unquam cogatur ad bonum aut cogi etiam possit. Si autem cogitur, in diversa rapitur) et coactionem sui odit. Sola autem gratia trahitur, hoc est, vere liberum efficitur, ut Augustus contra Julianum lib. 2).

Ubi Christum dixi vocasse apostolos et Magdalenam blando tractu, contra opponit, Paulum vocatum esse durius, et in evangelio iussum esse servum, ut compelleret intrare. Primum satis admiror, quod cum nos exemplum Pauli soleamus adducere pro defendenda gratia, ipsi nobis elabuntur dicentes, hoc esse miraculosum factum et non facere regulam. Hic tamen egregius D. D. quasi non pro miraculo habens, pro regula inducit. Sed hoc dimitto. Dico, non potuisse Paulum intus ex corde converti, nisi trahente gratia, ut b. Augustus contra epistolas Pelagianorum docet, gratiam dei facere ex nolentibus ac resistentibus volentes et sequaces. Ita et illud in evangelio solvitur, quod servus potest compellere verbo, sed nisi Dominus sibilet, nihil sequitur).

Ad illud Augustini, intelligentis servilem timorem in verbo Salomonis: Initium sapientiae timor Domini, et illud glossae Psal. 111: Metus iudicii ianua est conversionis ad deum, respondeo: Si adsit)

gratia; alioquin servilis timor sine gratia (neque enim Augustus excludit) non operatur nisi iram. Oportet enim, ut) Augustus non pugnet, sicut vere non pugnat, cum Paulo damnante omnia, quae sunt extra gratiam. Glossam super illud) Roma. 8 [v. 15]: Non accepistis spiritum servitutis in timore, dicentem, esse timorem servilem) bonum, potius reiicerem, tamquam expresse contra textum loquentem. Apostolus enim dicit: Non accepistis spiritum servitutis in timore, damnans eum. Aut dico, quod glossa) textum non exponit.

Illud apostoli ad Timoth. [1. 5 v. 20]: Argue coram omnibus, ut et) ceteri timorem habeant, D. D. ad servilem timorem adaptavit, quod relinquo. Ego de filiali timore accipio, donec aliter probet. Quod autem Augustus magistro reci-\marginpar{\small[371]}tante lib. 3. dicit: Timor servilis est, cum per timorem gehennae) continet se a peccato, quo praesentiam iudicis et poenas metuit etc.): continet, inquam, se a peccato externo tantum, intus tamen auget odium iustitiae gehennam comminantis. Et infra: Bonus est timor et utilis, licet) insufficiens, per quem paulatim fit consuetudo iustitiae. Hoc est meo iudicio consuetudo desperandi et odiendi deum, si excludatur gratia: verum autem est, si includatur gratia.

Non est ergo) necesse, ut dem aliam solutionem, nisi D. D. egregius primo) probaverit, s. patres loqui de servili timore, exclusa gratia, aut charitatem non habere negotium in expellendo timore servili. Quare, quae consequenter induxit de expulsione timoris, de initio sapientiae et quomodo expellat charitas timorem successu per augmentationem gratiae, ex praecedentibus satis intelliguntur).

Superest robustissima auctoritas Matth. 10 [v. 28.], quam unam sufficere putat: Nolite timere eos, qui occidunt corpus, sed timete eum etc.), ideo non esse servilem timorem damnabilem, ad quem Christus nos invitat. Respondeo primum: Si hoc vult D. D., contradicetur praecedentibus, ubi timor servilis insufficiens dicitur, ideo aeque absurdum) est dicere, Christum nos docuisse insufficientia). Dico tamen ego, non esse ibi servilem timorem 'doctum'), quia et filialis timor timet offendere deum et separari ab eo. Atque etiamsi de mero servili timore loqueretur, nondum intelligitur exclusa gratia, imo inclusa, cum teste apostolo et Augustino omnis lex et doctrina sit litera requirens spiritum.

Consensit etiam scholasticos vere dicere: Contritionem non valere extra charitatem, sed non confutavit. Stat ergo eorum ipsorum auctoritas contra eos: nisi confutet contritionem in charitate esse id, quod a charitate incipit fieri).

Illud Augustini bene placet: quod timor est medicamentum, charitas est sanitas, timor scilicet in charitate imperfecta, et charitas perfecta.

Rationem illam sumptam ex Dionysio), de imo, medio et summo, dicit me transiliisse) et non aperuisse, quod sit medium inter peccatum imum et charitatem superiorem. Fateor me libenter transiliisse), ut qui crederem ipsi D. D. satis esse perspectum), hanc rationem nihil esse ad propositum. Dionysius enim loquitur de ordinibus infimis, mediis, summis). At ego prorsus nullum habeo medium inter peccatum et gratiam, sicut nec Christus, quando dicit: Qui non est mecum, contra me est. Et iterum: Aut facite arborem bonam aut facite arborem malam. Sed et ipsum D. D. credo id ipsum asserere, quod gratia et peccatum apud scholasticos immediate opponuntur.

Non accipit D. D., quod gratiam dixi praevenire amorem et timorem, nisi loquar) de gratia primae motionis. Respondeo: Ego cum apostolo et Augustino sentio: quod nisi lex diligatur) (quod est charitatis et non primae motionis), non bene vivitur. Ergo nec deus timetur seu) colitur.

\marginpar{\small[372]} In fine instruit adversus me cornutum syllogismum: aut Ambrosius, Gregorius aliique repugnant dictis apostoli vel non; si sic, ut) ostendam; si non, ut stem eorum sententiae. Respondeo et per medium transeo: non repugnant dictis apostoli, si) et sto eorum sententiae, non autem Eccianae intelligentiae, imo errori. Non enim excludunt charitatem a timore sive servili sive filiali.

Dicit D. D. esse) cancrino more retrogradi, quod dixi: Habita charitate moveri mentem ad timorem. Satis miror sive scorpionem, sive cancrum illum), cum et gentilis ille poeta dixerit): Res est solliciti plena timoris amor. Quasi ignoremus, quod amor fons et caput est omnium affectuum. Ideo enim timetur poena et infernus serviliter, quia diligitur vita et voluptas pueriliter. Et ad hunc amorem et timorem) extirpandum diffunditur charitas dei, qua amemus aliam vitam et timeamus aliam mortem, id est, separationem a deo.

Transiit D. D. de punctis in sermone meo signatis et de toto

sermone discutiendo. Dico breviter: Habet calamum et papyrum, aggrediatur eum cum fiducia, videbitur, an sermonem confutaverit, an se ipsum irriserit. Haec relinquo iudicio ordinandorum.

\paragraph*{D. D. Eccius.}

Quia 'ob' temporis angustiam respondere non possum his, quae reverendus pater adduxit: refero tamen me ad iudices, me nihil transiliisse neque fuisse de mente praedicatorum aut doctorum, legem adimpleri sine charitate, neque aliquos docuisse infructuosam poenitentiam, sed quomodo per timorem servilem ad fructuosam perveniatur. Et quod magnifico verborum apparatu recessit ab hodierna solutione, [D. D. Martinus dixit: non recessi] qua dixerat Sapientem loqui de timore filiali, iam vero admittit loquentem de timore servili: non tamen excludendo gratiam, quod nec textus, nec s. doctores patiuntur. Nam per sapientiam intelligit charitatem, et sic initium esset ante initium. Totus quoque Augustinus non laborat, quomodo timor servilis cum charitate sit initium charitatis, sed quomodo timor servilis primo occupet mentem et primo ingrediatur et sic inducat charitatem, quod etiam est de mente Gregorii improbantis reditum cancrinum, ubi reverendus pater pro Gregorio scorpionem mihi adducit: Res est solliciti etc.. Omnes ergo s. doctores hodie citati volunt, timorem servilem esse initium charitatis, ad intelligentiam saepe datam, et quam praedicatores solent dare, de quo me refero ad iudices.

Crastina hora secunda propter festum divae Margaritae disputabunt iidem D. doctores in materia absolutionis a poena et a culpa per quemlibet simplicem sacerdotem.

\marginpar{\small[378]} Die Mercurii tertia decima mensis Julii hora secunda pomeridiana.

\paragraph*{D. D. Eccius.}

Circa materiam conclusionum quartae et quintae, reverende pater, quia tempus nobis praefixum est nimis artum, tamen tangendo fundum

negotii aliqua obiter attingemus, illud scilicet, quod quilibet sacerdos absolvat a poena et culpa, contra communem usum totius ecclesiae, et quod in sermone vulgari et conclusione dicitis, ex nulla scriptura probari, divinam iustitiam aliquam poenam vel satisfactionem expostulare a peccatore. Et illis innitor, quia in absolutione sacramentali culpa remittitur simultanee, sed poena peccato debita aeterna commutatur in temporalem. Hoc liquet ex sacrae scripturae traditione et s. patrum usu. Et ut brevior sim, expresse hoc voluit Ambrosius super Lucam, Hieronymus lib. 1. contra Jovinianum, Augustinus quaesti. 1. de octo Dulcitii quaestionibus, et pertinenter Ambrosius Luc. 5. testatur satisfactione poenam peccati dissolvi.

Scriptura est in promptu: quoniam peccatum Adae, culpa etiam remissa, punitur in tota posteritate. Quare transeunte culpa remanet poena. Quam rationem tangit Augustinus tractatu 124. super Matth. et glossa 2. Reg. [Sam.] 12. Similiter de David. 2. Reg. 24 [v. 10], ubi David poenituerat et percusserat cor eius: peccavi, dixit, valde in hoc facto, sed peto, transferas iniquitatem servi tui. Modo peccatum non transfertur quoad culpam, quia, ut propheta ait, anima ipsa quae peccaverit morietur. Restat ergo, ut peccatum quoad poenam transferatur. Sic per prophetam Nathan dicitur David: Transtulit deus peccatum tuum. Ubi glossa: Deus delictum delet, sed inultum non deserit. Aut enim homo in se poenitens punit, aut deus cum homine vindicans percutit.

Et ut ad s. patres veniamus, a quibus usus et praxis ecclesiae principaliter ab ipsis apostolis derivatus est. Nam primo nullum peccatum manet impunitum, et poena est ordinativa culpae. Quas rationes tangit Augustinus et refert Gratianus can. Sicut primi de poenitentia, distinctione 1.. Nam deus malum fieri non pateretur, nisi illud per iustitiam ordinando melius stare in mundo faceret. Inquit ergo Augustinus lib. de poenitentiae medicina: Non sufficit mores in melius commutare et a factis malis recedere, nisi etiam de his, quae facta sunt, satisfiat deo per poenitentiae dolorem, per humilitatis gemitum, per contriti cordis sacrificium, eleemosynis cooperantibus etc. Non enim dictum est, ut tantum abstineatis a peccatis, sed et de praeteritis, inquit, deprecare Dominum, ut tibi dimittantur etc.

\marginpar{\small[374]} Sed reverendus pater, sicut in resolutorio et alias, meo quidem sensu, facere nititur, satisfactionem illam fieri ecclesiae, sed deum

[kalligraphisch geschrieben]; de absolutione sacerdotis (sacerdotali?); a. R. circa conclusionem 4 et 5; U Tertia decima Julii Mcccc. XIX. hora secunda. Circa materiam conclusionum quartae et quintae; a. R. Eccius.

talem poenam non requirere, et quam deus requirat, homo non auferat; contra ego oppono: Primo quod Augustinus in Enchiridio inquit, cap. 71. (loquitur de oratione dominica): Delet omnino haec oratio minima et quotidiana peccata, delet et illa, a quibus vita fidelium scelerate gesta, sed poenitendo in melius mutata discedit. Ubi constat, orationem dominicam delere venialia et etiam mortalia, non quoad culpam, sed quoad poenam, quia poenitendo mutata, quae mutatio Augustini intelligi nequit, nisi de mutatione poenae aeternae in temporalem.

Eos, qui fuerunt a quadringentis annis, non adduco, ut Wilhelmum Parisiensem, Wilhemum Altisiod., qui id pleno ore affirmant. Accedat Cyprianus in epistola ad Fidum, ubi sic ait: Legimus literas tuas, carissime frater, quibus significasti de quodam presbytero Victore, quod ei, antequam plenam poenitentiam egisset et Domino deo, in quem deliquerat, satisfecisset, temere Therapius collega praepropere pacem dederit. Hic s. Cyprianus non dicit Victorem ecclesiae non satisfecisse, sed Domino deo.

Idem Cyprianus ad Cornelium papam de haereticis ait: Elaborant, ut opus suum diaboli malitia consumment, ne vulneratos divina clementia in ecclesia sua curet. Miserorum poenitentiam mendaciorum fraude corrumpunt, ne deo indignanti satisfiat. Apertissime sanctus martyr haereticos increpat, qui poenitentiam ab ecclesia iniunctam existimabant non esse satisfactionem dei.

Et quamvis reverendus pater citra modestiam saepius mihi improperet, quasi scripturarum sanctarum intelligentiam non idonee tractem, de quo iudicent hodie designati iudices, uter rectius de fide et sacrarum scripturarum sensu sentiat: attamen, quantum ego video, universi sacrarum scripturarum tractatores in hanc sententiam descendunt. Ad quid enim ecclesia vellet imponere tale opus grave poenitentibus, cum Dominus et caput ecclesiae illud non exigeret? Quam rationem in annotationibus tetigi, non tamen fuit mea, sed Chrysostomi, sicut refert Gratianus 26. quaest., 7. can. Alligant: homo cui grave pondus poenitentiae imponis, aut poenitentiam reiiciet, aut dum sufferre nequit, scandalizatus amplius peccat. Deinde etsi erramus, poenitentiam modicam imponentes, nonne melius est, propter misericordiam rationem reddere quam propter crudelitatem? Ubi enim paterfamilias largus est, dispensator non debet esse tenax. Si deus benignus, ut quid sacerdos vult esse austerus? Hoc pro confessoribus.

Idem Homilia 31. ait super epistola ad Hebraeos: Peccatum confessione minuitur, sed non minuitur quoad culpam, quia impium est, ut Hieronymus ait, a deo dimidiam sperare veniam, ergo minuitur quoad poenam.

\marginpar{\small[375]} Addo Gregorium, 4. moral. § 49., ubi post multa ita dicit: Sed quia nullum peccatum deus inultum relaxat, aut enim nos hoc flendo insequimur, aut ipse iudicando, restat, ut ad emendationem suam semper mens solerter invigilet etc..

Hoc sensit et b. Hieronymus, verba sua transcripta sunt can. Mensuram, de poenitent. dis. 1: Mensuram temporis in agenda poenitentia idcirco non satis praefigunt canones pro unoquoque crimine, ut de singulis dicant, qualiter unumquodque emendandum sit, sed magis in arbitrio sacerdotis intelligentis relinquendum statuunt, quia apud deum non tam valet mensura temporis quam doloris.

Augustinum quoque refert in can. Nullus: Nullus debitae gravioris poenae accipit veniam, nisi qualemcunque etsi longe minorem, quam debeat, poenam solverit, ita enim a deo impertitur largitas misericordiae, ut non relinquatur iustitiae disciplina. Hic habet egregius D. D. Carolostadius Augustinum in defensione mea adductum, cuius sententiam prosequi non est instituti nostri in praesentia.

Accedat postremo Isidorus lib. 2. de summo bono, cap. 13: Quamvis per poenitentiam propitiatio peccatorum sit, tamen sine metu hoc esse non debet, quia poenitentis satisfactio divino tantum pensatur iudicio, non humano. Quare poenitentia iniuncta non est, quia satisfiat tantum ecclesiae, sed etiam quia satisfit deo. Nam et apud deum productior est poena quam culpa, ut Augustinus testatur tract. 124 super Iohannem.

Quibus omnibus pro indubitata veritate habitis faciliter patet improbatio conclusionis quintae, ubi dicit reverendus pater, quemlibet sacerdotem debere absolvere a poena et a culpa, et ita quilibet sacerdos villanus esset in sua parochia episcopus, archiepiscopus et papa. Quod esse manifestissime falsum praeter usum totius ecclesiae, non solum adulatorum, probatur primo ex his, quae dicta sunt, quia per absolutionem sacerdotalem culpa deletur, poena manet, licet commutata; deinde, quia sacramentum poenitentiae est iudiciale, ut formam huius iudicii describit b. Augustinus ca. forte 2. de poen. medic. Et ad iudicium ferendum pertinet iurisdictio, nisi velit Anaxagoreum chaos et confusionem maximam facere in ecclesia dei.

Et cum iurisdictio se ad duo extendat materialia in hoc sacramento habita, nam et peccator poenitens et peccata ipsa sunt de materia poenitentiae et nemo negat iurisdictionem in inferioribus praelatis ad confusionem tollendam esse coartatam quoad materiam peccantium, nam hinc proprii sacerdotis habetur determinatio, et quia sententia non lata a suo iudice est nulla, ideo absolvens non subditum nihil facit.

Sic a simili in materia peccatorum potest coar- \marginpar{\small[376]} tari iurisdictio eadem ratione et in criminum detestationem. Quod autem etiam a sibi commissis non possit plenarie absolvere, ex eo patet, quia tanta potestas, nisi auctoritate fundetur, contra totius ecclesiae consensum, cuius praxis debet homini Christiano esse pro regula, non debet tribui. Sed in superioribus ostensum fuit per maiores clericos, Cyprianum, Chrysostomum et Augustinum, poenae debitum remanere post culpam. Ideo reverendus pater cedat s. patrum auctoritati et consuetudini totius ecclesiae obtemperet, aut cur id fieri non debeat, auctoritate sacrae scripturae, conciliorum vel sanctorum patrum exponat.

Dixit in fine D. D. Eccius, se nihilominus velle replicare, si necessarium videretur sibi, post responsionem domini doctoris Martini.

D. D. Martinus respondit.

Admitto D. doctori, quod habeat ultimum verbum, quia sic vult. Arguit autem egregius D. D. contra conclusionem meam, volens probare poenam requiri a deo, et pro primo inducit scripturam Gen. 3, ubi peccatum Adae etiam remissa culpa punitur in posteritate, ergo poena remanet, transeunte culpa. Volo D. doctorem obstrictum hoc exemplo, et non replicet, nisi confirmet, aut sententia cedat. Si ista poena requiritur pro peccato et similiter sentiendum est de qualibet poena, ut infert, habeo propositum, quod poena non possit remitti per papam aut ullum sacerdotem, quia nullus adhuc remisit mortem, infinita genera poenarum, morborum et similium miseriarum, quae omnes pro peccato primo sunt illatae. Et sic probatio D. doctoris probabit conclusionem meam contra se ipsum.

Secundo, quod a David translato peccato non est translata tamen poena, etiam pro me facit, sicut et in sermone eodem vulgari sequenter dixi: Quia erat poena, quam deus requisivit, ideo nemo

auferre potest. Si D. D. non intellexit me, legat diligentius, ego enim dixi et dico, quod hanc poenam deus non requirat, quam papa vel homo possit solvere; quam autem ipse requirat, scilicet immediate, non possit homo solvere.

Tertio, induxit glossam: Deus delictum delet, sed inultum non dimittit; nihil pugnat contra me. Ulciscitur enim, sive per ipsummet hominem conterendo, sive per ecclesiam emendando, sive per se ipsum iudicando. Et hoc ultimum et primum genus poenae in nullius hominis arbitrio est, sicut apostolus 1. Cor. 11 [v. 31 f.] dicit: Si nos ipsos iudicaremus, non utique iudicaremur a Domino, cum autem iudicamur a Domino, corripimur, ut non cum hoc mundo damnemur. Ibi clarum est, quomodo deus poenam requirat et non requirat.

Quarto, Augustinum per Gratianum relatum, quod deus non permitteret malum fieri, nisi per iustitiam illud ordinando melius stare faceret: miror egregium D. D.: quod has et similes auctoritates ita consarcinat, cum nemo eas neget aut contrarium sapiat. Ego enim hoc solum impugnavi semper, quod virtute \marginpar{\small[377]} clavium iactant solvi poenas per iustitiam divinam requisitas. Hoc enim non credo esse verum, nec probabitur. Non enim dixit ad Petrum: Quod ego ligo, tu solves, sed: Quodcunque tu solveris, solutum erit [Matth. 16 v. 19].

Quinto, illud de poenitentiae medicina: Non sufficit mores in melius commutare, nisi etiam satisfiat deo de his, quae facta sunt, per poenitentiae dolorem, per humilitatis gemitum, per contriti cordis sacrificium. Haec omnia ego quam maxime volui semper, ut quid ergo per indulgentias ista iactantur remitti, si non sufficit mores in melius commutare et divina iustitia has requirit satisfactiones? Frustra ergo gloriatur D. D., quod non dictum sit, ut tantum abstineatis a peccatis, sed de praeteritis deprecare Dominum. Et sic patet, quod adhuc nulla scriptura contra me adducta est, quanquam ego possem multo fortiores inducere pro me.

Post haec adducit Augustinum in Enchiridio: Delet omnino haec oratio minima et quotidiana peccata, delet et illa, a quibus vita fidelium scelerate gesta, sed poenitendo in melius mutata discedit. Hoc est, quod

dixi, peccatorem post conversionem teneri ad crucem et passiones vitae, secundum quod deus intulerit, quas non possit homo solvere, et iterum haec auctoritas pro me facit.

Eos, qui fuerunt a quadringentis annis, non adducit, et placet.

Ad Cyprianum, in epistola ad Fidum, ubi damnat Therapium, quod praepropere pacem dederit Victori, antequam plenam egisset poenitentiam et Domino satisfecisset: Ecce, inquit, non ecclesiae, sed Domino 'non' satisfecisse dicit Victorem. Respondeo: Legat et conferat Cyprianum bene D. D. et inveniet eos etiam, quos paci dederant, ideo datos esse paci etiam praepropere, ut cruces et martyria expeditius sustinerent. Quae ipse per multas epistolas exponit esse poenas et flagella peccatis a deo inflicta. Quare Victor hoc modo nondum satisfecit deo; et tamen satisfecit, quia ecclesiae, quam nos vult audire deus, satisfecit. Hoc enim habent verba Christi, quod ecclesia debet imponere poenas, quando dicit: Quodcunque ligaveris. Et hoc modo possem admittere, deum requirere poenas, quas ecclesia possit solvere, quia pactum fecit cum illa.

Altera auctoritas Cypriani ad Cornelium, qua iterum scribit, haereticos impedire peccantes, ne indignanti deo satisfaciant: Quanquam Cyprianus ibi de poenitentibus non loquatur, sed de haereticis, qui tanquam iusti et quasi bene fecissent excusaverunt et defenderunt se in peccatis suis, tamen respondeo sicut ad priora.

Post haec addidit rationem: Ad quid ecclesia vellet onus imponere poenitentibus, si illud Dominus non exigeret, et longam post auctoritatem ca. Alligant inducit, quae omnia transeo, et dico, quod ecclesiae mandatum est, quod castiget et iudicet peccatores, quod si non fecerit, deus irremissibiliter faciet iuxta sententiam Pauli superius inductam 1. Cor. 11 (v. 31 f.) et sic non potest solvi.

\marginpar{\small[378]} Item idem de Homilia 31. super Hebraeos: Peccatum confessione minuitur. Ex quo verbo colligit D. D. in hunc modum: Non minuitur quoad culpam, quod impium sit a deo sperare dimidiam veniam: ergo quoad poenam. Admitto totum, secundum praedicta.

Iam illa auctoritas Gregorii 4. Mor.: Aut enim nos hoc flendo insequimur, aut ipse iudicando; vix aliud aptius pro me adduci potuit. Similiter et illud Hiero. sequens c. Mensuram, quod tamen Hieronymi esse dubito, admitto totum, quod canones non satis praefigunt, ideo relinqui arbitrio sacerdotis. Addo et ego: multo magis arbitrio dei, qui solus est ponderator spirituum et non ignarus neque iniquus iudicat.

Placet et illa auctoritas Augustini in cap. Nullus: Ita impartitur a deo largitas misericordiae, ut non relinquatur iustitiae disciplina. Potest pro me et pro D. D. valere.

Et Isidorus quamquam non satis gravis auctor in his rebus, placet tamen, quod poenitentis satisfactionem tantum divino pensari iudicio dictat, non humano. Ergo multo minus remitti potest humano iudicio, cum clavis potestatis non debeat operari, nisi prior sit clavis scientiae, quae sciat, quid et quantum solvat.

Apud deum productiorem esse poenam quam culpam, ex Augustino concedo salva tamen auctoritate apostoli Roman. 7. [v. 18] dicentis: Non invenio in me, hoc est, in carne mea, bonum. Cuius sententia est, poenam et peccatum simul finiri.

Haec de prima conclusione. — Contra aliam dicit, esse manifestissime falsum et praeter usum totius ecclesiae, quod quilibet sacerdos debeat absolvere a poena et culpa poenitentem. Et hoc probat primo ex dictis, id est ex nihilo. Secundo per rationem, quia sacramentum poenitentiae sit quoddam iudiciale, et ad iudicium ferendum pertinet iurisdictio. Iurisdictionem autem esse coartatam in inferioribus praelatis ad tollendam confusionem, tam in peccantibus quam in peccatis, secundo in detestationem criminum, alioquin quilibet villanus sacerdos esset episcopus, archiepiscopus et papa.

Respondeo et dico duo, primo me nescire usque in hodiernum diem, an coartatio istius iurisdictionis fecerit hoc, quod praetenditur, scilicet criminum detestationem et confusionis ablationem. Hoc certum est, longe secius evenisse, nam peccata crassissima etiam ridentur in maioribus curiis, quae in propriis parochiis possent pulcherrime puniri, si servatus esset modus, quem instituerunt apostoli et servaverunt s. patres usque post Nicenum concilium, ubi definitum est et longe post servatum, ut dioeceses non permiscerentur et unusquisque poeniteret in sua dioecesi. De \marginpar{\small[379]} quo exstat cum aliis praeclarissima epistola Cypriani, ad Cornelium Romanum pontificem epistola III: Nam cum statutum sit ab omnibus nobis et aequum sit pariter et iustum, ut uniuscuiusque causa illic audiatur, ubi est crimen admissum, et singulis pastoribus portio gregis sit adscripta, quam regat unusquisque et gubernet, rationem sui actus Domino redditurus, oportet utique eos, quibus praesumus, non circumcursare nec episcoporum concordiam cohaerentem sua subdola et fallaci temeritate collidere, sed agere illic causam suam, ubi et accusatores haberi et testes sui criminis possint. Loquitur enim de his, qui in Africa peccaverant et ad Romanum pontificem Cornelium cucurrerant. Et sic patet usus primitivae ecclesiae de ligandis et solvendis peccatoribus.

Cum autem, ut apostolus Act. 20. [v. 17 u. 28] ostendit, idem sit quod episcopus et presbyter, et ad Tit. 1. [v. 5] quaelibet civitas suum episcopum habere debeat iure divino, longe utilior esset ad corripienda peccata modus, si quilibet sacerdos in sua parochia ligaret et solveret poenitentem. Quod exemplum monstravit apostolus 1. Corinth. 5, ubi cum Corinthiis praesens spiritu tradidit fornicarium satanae, obiurgans quod ipsi non fecissent.

Quo vero iure aut qua ecclesiae felicitate hic modus divino iure praescriptus et in tantum tempus roboratus sit sublatus, viderint alii. Ego quidem negare non possum, quia ad oculum videmus ita fieri, ita statui, ut tam peccantes quam peccata reserventur et uni animae sex vel septem pastores secundum, sub et supra imponantur. Sed an ita debeat fieri aut expediat, non definio. Scio, quod inferior tenetur obedire restringenti et vexanti, quanquam nullo iure divino. Superior tamen, ut dicit conclusio mea, gravissime peccat, si reservet occulta peccata sine rationabilissima causa, imo adhuc dubito, et quantum capio, credo sine temeritate, nullum occultum peccatum debere reservari aut posse, optans audire probationem contrariam.

Secundo dico, ecclesia non esset ruitura, si idem plebanus episcopus, archiepiscopus et papa esset ac sola concordia cohaerente, ut Cyprianus ait, et sicut usus prioris ecclesiae fuit, iungerentur.

Proinde, quod et damnatus fertur articulus in Constantiensi concilio de istis reservationibus, non satis curo. Hoc scio, quod probatus fuit et reprobata ista reservatio in primitiva ecclesia et institutione apostolorum, et nunc quoque, ut miserrima ecclesiae experientia docet, esset utilissimus et saluberrimus ad coërcenda peccata et tollendam detestabilem confusionem omnium episcopatuum, quam hodie videmus. Qua causa autem sit suo tempore mutatus, transeo. Mutabilem invenio, mutabilem relinquo. Haec in arbitrium iudicum.

\marginpar{\small[380]}
\paragraph*{D. D. Eccius.}

Reverendus pater ab initio satis glorianter parvi pendit per me adducta, quasi sibi minime obstent, quae maxime adversantur. Nam ipse voluit in sermone vulgari, non fieri commutationem poenae aeternae in temporalem, et opinioni communi contradicit, quasi non possit probari, deum exigere aliquam satisfactionem, praeter portationem crucis. Et in sermone latino collaudat

Haec coepta fuerunt tantum per D. D. Eccium circa horam quartam vespertinam, sed intercepta et reassumpta in exordio disputationis postridianae, videlicet

hora septima antemeridiana die Jovis XIV. mensis Julii.

dictum vulgarium ultra omnem doctrinam doctorum scholasticorum, de poenitentia datam: NVMMER THVN DIE HOCHSTE BVES. Optima poenitentia nova vita, secundum glossam eius. Addiderunt talem persuasionem ipse et propugnator: si averterit se impius etc.. Si ad novam poenam imputat, quomodo dicitur non recordari? et patet conclusio sua manifeste dicens, peccare sacerdotem, qui non absolvat a culpa et poena. Huic errori ego contradixi s. patrum auctoritatibus, quibus ipse voluit illudere et auditores fascinare, quasi apud eum solum esset potestas interpretandi sacram scripturam.

Unde bene adduxi contra eum Augustinum: Non sufficit mores in melius commutare et a malis factis recedere. Ubi est liquidissimum, doctrinam ex sermone allegatam esse falsam, quia nova vita non est optima poenitentia, cum non sufficiat secundum Augustinum. Et bene etiam illa doctrina improbatur per b. Ambrosium, lib. 2. de poeniten. cap. 5: Apostoli secundum Christi magisterium docuerunt poenitentiam. Et infra: Qui enim agit poenitentiam, non solum diluere debet lacrimis peccatum suum, sed etiam emendatioribus factis operire et tegere peccata sua. Clarissima s. patris verba, quod emendatioribus factis agendo poenitentiam tegere debemus peccata.

At cum auctoritates forent expressissimae, deum non remittere peccatum impunitum, confugit ad mirabilem poenarum distinctionem, qui tamen solitus est scholasticos ob distinctionum usum improbare et dixit poenam, qua deus vult peccatum punire, non posse auferri per hominem vel per papam, quod est omnium falsissimum et potestatis clavium annullativum. Nam in poenitentia facta commutatione poenae certe homo potest solvere illam poenam, per ipsum apostolum, per reverendum patrem inductum, 1. Cor. 11: Si nos ipsos iudicaremus, non utique iudicaremur a Domino. Quare si nos pro hac poena satisfacimus, deus pro peccato a nobis aliam non exigit; alioquin contra prophetam, illud quod nos puniremus, si deus vellet punire, iam bis puniret id ipsum.

\marginpar{\small[381]} Praeterea clara fuerunt Cypriani verba, Chrysostomi, Gregorii et Hieronymi, quod poena iniuncta a sacerdote in satisfactione est poena deo debita, et ita Victor, quia non impleverat poenitentiam iniunctam, a Cypriano dicitur nondum satisfecisse deo. Quod et Theodorus in suo poenitentiali, quem sequitur quasi ad verbum Beda, id ipsum testatur, quamvis forte nec ii auctores sicut nec Isidorus ei satisfaciunt.

Unde quidquid dicat ore, auctoritates omnes facere pro eo, tamen necessarium est, quod corde dissentiat: cum hae auctoritates doctrinam eius et conclusionem subvertant. Si enim sacerdos peccat non absolvendo a poena et culpa, tunc episcopi peccassent non absolventes Victorem a poena et culpa, et omnes sacerdotes peccarent per orbem Christianum non absolventes extra indulgentias.

Neque relevatur in eo, quod debeat crucem portare et hanc poenam exigat deus, quia ista crucis portatio non est aliud, quam vita Christiniana, sicut ipse reverendus pater docte hoc exposuit; sed cum hoc oportet de praeteritis satisfacere et Dominum deprecari de praeteritis.

De reservatione casuum conclusio mea nihil habet, tamen reservationes moderatas credo esse utiles, quod etiam praelati in monasteriis experiuntur. Fateor ingenue, me esse in voto Gersonis, in concilio Constantiensi interpellantis pro delenda immodica reservatione casuum, et potissimum displicet mihi illa reservatio, sicut et domino patri, quando habet avaritiam comitem, id est, pecuniariam poenam annexam.

Ecclesiam non dicit ruituram, si sacerdos in sua parochia esset episcopus et papa; certe apparet mihi, si esset illius pulcherrimi ordinis hierarchici corruptio, quod etiam esset ecclesiae subversio.

Sed maneam in principali, quia peccatum non manet impunitum, secundum Augustinum, Gregorium etc.. Quare satisfactio merito tertia pars poenitentiae asseritur. Et plene probavit Augustinus in Ench., per orationem dominicam nos satisfacere pro scelerate gestis. Et ut vere dixit heri reverendus pater, deus fecit pactum cum ecclesia: Si ecclesia non facit, deus facit. Ergo habeo propositum scholasticorum et praedicantium, quod oportet vel nos satisfacere, aut deus exiget. Quodsi per orationes nostras vel bona opera satisfacimus, cur illud non posset fieri accedente virtute clavium, quas deus non frustra suae sponsae ecclesiae contulit? Et cum iuxta Gregorii sententiam cap. Decreto 2, q. 6, alii episcopi sunt vocati in partem sollicitudinis, papa habet plenitudinem potestatis, per indulgentias ab eo datas satisfit poenae pro peccatis debitae solutione ex thesauro ecclesiae facta, ut post Sixtum declaravit papa modernus, praecipiendo sub poena excommunicationis ita doceri, \marginpar{\small[382]} teneri et praedicari. Quare si reverendus pater contra doceat, praedicet vel disputet, iam est anathemate percussus. Placuit tamen mihi, quod ultra clavem potestatis heri posuit clavem discretionis, cum tamen plures claves contra scholasticos doctores neget conclusione 7. Resolutorii.

Haec volui omnino adiicere, ut futuri iudices apertius intelligerent, quae nostra esset controversia in hoc puncto, quoniam, si reverendus pater stat sententiae doctorum per me allegatorum, non adversabitur nec scholasticis nec praedicatoribus nec mihi. Quare si voluerit, etiam

sententiam suam pro informatione dominorum iudicum poterit clarius exprimere.

D. D. Martinus
non ascensa cathedra, sed ex loco sedendi solito
assurgens respondit.

Ad istas naenias et ineptias D. doctoris heri satis respondi, repetit enim eadem et velut ridiculus citharoedus chorda semper oberrat eadem.

Secundo scopum controversiae non attigit Non enim quaestio est, an deus peccatum inultum dimittat, quod satis copiose probavit, sed an papa vel ecclesia remittat, quas poenas deus exigat. De hoc nihil probavit, quod relinquo iudicio iudicum et omnium auditorum.

Tertio, obticuit hodie scripturas sanctas, ideo sto in eius probatione hesterna prima, Gene. 3. ca., ubi probavit, poenas requiri a deo, quas ibi scriptura ostendit irremissibiles. Doleo, quod D. D. ita profunde penetrat scripturas sicut tipula aquas, imo videtur fugere a facie earum sicut diabolus crucem. Quare salvis reverentiis patrum, praefero ego auctoritatem scripturae, quod commendo iudicibus futuris.

\paragraph*{D. D. Eccius.}

Quia impatiens monachus scurrilia quaedam addidit praeter gravitatem theologicam, de quo integri viri iudicent, an recte induxerim contra eum, iudices iudicabunt. Sed hanc fuisse materiam nostram, patet ex conclusione 4.: dicere, deum remittendo culpam remittere poenam etc. Hoc erat saxum volvendum, et quia praefert auctoritatem sacrae scripturae patribus, quasi ipse velut alterum oraculum Apollinis solus habeat scripturarum intelligentiam ultra s. patres, et apparenter inducit auctoritatem heri per me citatam, dico duo.

Primo me adduxisse in eum finem, ut ostenderetur, cum impius avertisset se ab iniustitia sua, deum adhuc recordari, non quoad culpam, sed quoad poenam, quod facit pro scholasticis et praedicantibus.

Secundo, cum reverendus pater reflectit auctoritatem illam contra me, quia poenae istae a deo peccato Adae impositae sint irremissibiles a papa et homine, verum est, et fateor, quia istae sunt poenae, quae

non solum consequuntur per-\marginpar{\small[383]}sonam, sed etiam consequuntur naturam, et ergo non est mirum, quod istae poenae non sunt remissibiles ab homine, sed per hoc non probatur, poenas personales pro peccato debitas non esse a sacerdote vel papa remissibiles. At in his refero me ad iudices, paratus mutare sententiam, si me aliter et melius docuerint.

